
\subsection{Proof of Theorem \ref{thm:approx-true-kernel}}\label{app:thm-approx-true-kernel}
\begin{proof}
    We split the proofs into several parts. First, we show that the characteristic function of the true distribution is sufficiently close to the one from the Gaussian distribution governed by $\bm \Gamma$. Then, we use that to bound $\|k_1^* - k_1\|_{L^2(\rho_{\mathcal X})\times L^2(\rho_{\mathcal X})}$, which immediately implies the operators are also close in norm. Since these objects live in the $L^2(\rho_{\mathcal X})$ space, the inputs are inherently unbounded, thus we further separate this case to deal with a bounded set and use the exponential decay of the measure to control the tails. Lastly, we show the distribution shift identity by introducing a pushforward measure dictated by $\bm \Gamma$ to translate the weights and inputs to the context of isotropic weights.
    \item \paragraph{Convergence of the true distribution to the spiked Gaussian:}
    Unconditionally, $\bm z$ is not Gaussian in general (it is a mixture of sub-exponential random variables), but conditional on $\bm \varsigma$ it is Gaussian.
    To be specific, if we let $c \coloneqq \frac{\mu_1 \eta}{n \sqrt{m}}$, with $a \sim \mathcal N(0, \frac{1}{m})$ and $\bm w \sim \mathcal N(0, \frac{1}{d}\bm I_d)$, condition on $\bm \varsigma = \sum y_i \bm x_i$, we have
    \[
        \bm z \ | \ \bm \varsigma \sim \mathcal N(0, \bm \Sigma(\bm \varsigma)), \quad \bm \Sigma(\bm \varsigma) = \frac{1}{d}\bm I_d + \frac{c^2}{m}\bm \varsigma \bm \varsigma^\top\, .
    \]
    So conditioning on $\bm \varsigma$, $\bm z$ is Gaussian with an anisotropic rank-one spike along the direction $\bm \varsigma$. Therefore, if we fix a test vector $\bm t \in \mathbb R^d$ and consider the projection $\langle \bm z, \bm t\rangle$, conditionally on $\bm \varsigma$ we have
    \[
        \langle \bm z, \bm t\rangle \ | \ \bm \varsigma \sim \mathcal N\left(0, \frac{\|\bm t\|^2}{d} + \frac{c^2}{m}\langle \bm \varsigma, \bm t\rangle^2 \right)\, .
    \]
    And consequently, the unconditional characteristic function is given by
    \[
        \phi_{\langle \bm z, \bm t\rangle}(u) = \exp\left(- \frac{u^2}{2}\frac{\|\bm t\|^2}{d}\right) \mathbb E_{\bm \varsigma}\left[ \exp\left(- \frac{u^2}{2}\frac{c^2\langle \bm \varsigma, \bm t\rangle^2}{m}\right) \right]\, .
    \]
    
    Now, to investigate the concentration of this variable, we look at the difference
    \[
        \left| \mathbb E_{\bm \varsigma}\left[ \exp\left(- \frac{u^2}{2}\frac{c^2\langle \bm \varsigma, \bm t\rangle^2}{m}\right) \right] -  \exp\left(- \frac{u^2}{2}\frac{c^2\mathbb E_{\bm \varsigma}[\langle \bm \varsigma, \bm t\rangle^2]}{m}\right)\right|\, .
    \]
    
    If we define $X = \frac{c^2}{m}\langle \bm \varsigma, \bm t\rangle^2$, $\mu = \frac{c^2}{m}\mathbb E_{\bm \varsigma}[\langle \bm \varsigma, \bm t\rangle^2]$ and $\alpha = \frac{u^2}{2}$, we perform a second-order Taylor expansion of the function $f(x) = e^{-\alpha x}$ around $\mu$. By Taylor's theorem, there exists some $\xi$ between $X$ and $\mu$ such that
    \[
        e^{-\alpha X} = e^{-\alpha \mu} - \alpha e^{-\alpha \mu}(X - \mu) + \frac{\alpha^2 e^{-\alpha \xi}}{2}(X - \mu)^2
    \]

    Taking the expectation of both sides yields:$$\mathbb{E}[e^{-\alpha X}] = e^{-\alpha \mu} - \alpha e^{-\alpha \mu}\mathbb{E}[X - \mu] + \frac{\alpha^2}{2}\mathbb{E}[e^{-\alpha \xi}(X - \mu)^2]$$Because $\mu = \mathbb{E}[X]$, the first-order term cancels out exactly ($\mathbb{E}[X - \mu] = 0$), leaving:$$\mathbb{E}[e^{-\alpha X}] - e^{-\alpha \mu} = \frac{\alpha^2}{2}\mathbb{E}[e^{-\alpha \xi}(X - \mu)^2]$$

    Since $\alpha = u^2/2 > 0$ and $X \geq 0$ (as it is a scaled squared projection), it follows that $\xi \geq 0$ and therefore $e^{-\alpha \xi} \leq 1$. Taking the absolute value, we can bound the difference directly by the variance of $X$
    \[
        |\mathbb{E}[e^{-\alpha X}] - e^{-\alpha \mu}| \leq \frac{\alpha^2}{2}\mathbb{E}[(X - \mu)^2] = \frac{\alpha^2}{2}\text{Var}(X)
    \]
    Substituting our definitions back in, we have $\frac{\alpha^2}{2}\text{Var}(X) = \frac{u^4}{8}\text{Var}(X)$. Given that $\frac{c^2}{m} = \mathcal{O}(\frac{\eta^2}{d^4})$, the variance of the sub-exponential variable yields the following bound
    \begin{equation}\label{eq:char-fn-bound}
        |\mathbb{E}[e^{-\alpha X}] - e^{-\alpha \mu}| \leq \mathcal{O}\left(\frac{\eta^4}{d^5}\right)
    \end{equation}
    
    \item \paragraph{Bounding the kernel difference:}
    We let $p_{\bm \Gamma}$ be the density function of the Gaussian distribution $\mathcal N(0, \frac{1}{d}\bm \Gamma)$ and $p_1^* = \mathbb E_{\bm \varsigma}[p_{\bm I_d}(\bm z \ | \ \bm \varsigma)]$ be the true density function from the non-Gaussian distribution followed by $z$.
    
    Denote $G(\bm w) := \sigma(\langle \bm w, \bm x\rangle) \sigma(\langle \bm w, \bm x'\rangle)$, the true kernel after the deterministic update is given by
    \[
        k_1^*(\bm x, \bm x') = \int G(\bm w) p_1^*(\bm w) \mathrm d \bm v\, ,
    \]
    while the Gaussian kernel $k_1$ is
    \[
        k_1(\bm x, \bm x') = \int G(\bm w) p_{\bm \Gamma}(\bm w) \mathrm d \bm v\, .
    \]
    We define a 2-dimensional vector $\bm{v}$ representing the two projections
    \[
        \bm{v} = \begin{bmatrix} v_1 \\ v_2 \end{bmatrix} = \begin{bmatrix} \langle \bm{w}, \bm{x} \rangle \\ \langle \bm{w}, \bm{x}' \rangle \end{bmatrix} \in \mathbb{R}^2\, ,
    \]
    then we can write $G$ in terms of this 2D variable: $G(\bm{v}) = \sigma(v_1)\sigma(v_2)$. The error between both kernels is the difference in expectations over this 2D plane
    \[
        |k_1^*(\bm x, \bm x') - k_1(\bm x, \bm x')| = \left| \int_{\mathbb{R}^2} G(\bm{v}) p_1^*(\bm{v}) \mathrm d\bm{v} - \int_{\mathbb{R}^2} G(\bm{v}) p_{\bm \Gamma}(\bm{v}) \mathrm d\bm{v} \right|
    \]
    
    Taking the Fourier transforms, we map this into the 2D frequency domain. In this context, the frequency variable is $\bm{u} = (u_1, u_2)$ and the error is given by
    \[
        |k_1^*(\bm x, \bm x') - k_1(\bm x, \bm x')| = \frac{1}{(2\pi)^2} \left| \int_{\mathbb{R}^2} \hat{G}(\bm{u}) [\Phi_1^*(\bm{u}) - \Phi_{\bm \Gamma}(\bm{u})] \mathrm d\bm{u} \right|\, ,
    \]
    where $\Phi$ are the respective characteristic functions of the distributions. We can see that the inner product between $\bm u$ and $\bm v$ gives
    \[
        \langle \bm u, \bm v\rangle = u_1 \langle \bm w, \bm x\rangle + u_2 \langle \bm w, \bm x'\rangle = \langle \bm w, u_1 \bm x + u_2 \bm x' \rangle\, .
    \]
    Therefore, if we define our test vector as $\bm t_u = u_1 \bm x + u_2 \bm x' \in \mathbb R^d$, $\langle \bm u, \bm v\rangle = \langle \bm w, \bm t_u\rangle$, the characteristic function for $\bm u$ is given by
    \[
        \Phi(\bm u) = \mathbb E[\exp(i \langle \bm u, \bm v\rangle)] = \mathbb E[\exp(i \langle \bm w, \bm t_u \rangle)]\, .
    \]
    This implies the characteristic function acting on $\bm u$ is precisely the characteristic function of a 1D projection of $\bm w$, as considered in the bound from \cref{eq:char-fn-bound}. Thus, we can plug in the bound and factor the $d$-dependent estimate out of the integral
    \[
        |\Phi_1^*(\bm{u}) - \Phi_{\bm \Gamma}(\bm{u})| \leq \mathcal{O}\left(\frac{\eta^4\ln^3 d}{d^5}\right) \|\bm{u}\|^2 \exp\left(-\frac{1}{2}\bm{u}^T \bm \Sigma \bm{u}\right) \, ,
    \]
    where we write the covariance matrix $\bm \Sigma$ as
    \[
        \bm \Sigma = \frac{1}{d}\begin{bmatrix}
            \|\bm x\|^2 & \langle \bm x, \bm x' \rangle\\
            \langle \bm x, \bm x' \rangle & \|\bm x'\|^2
        \end{bmatrix}\, .
    \]
    
    We will prove the kernels are close in norm in the product space $L^2(\rho_{\mathcal X}) \times L^2(\rho_{\mathcal X})$, which leads to the same conclusion to the Hilbert-Schmidt norm and the operator norm of the difference of integral operators. For this we split this proof with a truncation argument considering a bounded domain and then use the Lipschitz property of $\sigma$ to ensure the tails decays exponentially. 
    Ultimately, we want analyze the integral
    \[
        \iint|k_1^*(\bm x, \bm x') - k_1(\bm x, \bm x')|^2 \mathrm d \rho_{\mathcal X}(\bm x)\mathrm d \rho_{\mathcal X}(\bm x')
    \]
    so we start controlling the difference inside a bounded set.
    \paragraph{Bound of the pointwise difference over a bounded set: }
    To avoid singularities near the origin, we redefine our bounded domain as the set
    \[
        B_{R,c} = \{\bm x \in \mathbb{R}^d : c\sqrt{d} \le \|\bm x\| < R\sqrt{d}\}
    \]
    and we consider the bounded product set $B_{R,c} \times B_{R,c}$.
    Then, since $\bm \Sigma$ is positive semidefinite and inside this set it is never degenerate, we have 
    \[
        \Trarg{\bm \Sigma} = \frac{\|\bm x\|^2 + \|\bm x'\|^2}{d} \leq 2R^2\, .
    \]
    Now, because both weight distributions are sub-exponential, using Lemma \ref{lemma:activ-prod}, we get that
    \[
        \sup_{\bm x, \bm x' \in B_{R, c} \times B_{R, c}}|k_1^*(\bm x, \bm x') - k_1(\bm x, \bm x')| = \mathcal O\left(\frac{\eta^4\ln^3 d}{d^5}\right)\, .
    \]

    \paragraph{Bounding the integral over the tails:}
    Since the integration over the bounded set is handled, we analyze the integral over $( B_{R, c} \times B_{R, c} )^c$. Because the activation function $\sigma$ is $L_{\sigma}$-Lipschitz, we have the inequality $|\sigma(t)| \le |\sigma(0)| + L_{\sigma}|t|$. Consequently, for both kernels, there exist absolute constants $C_1, C_2 > 0$ such that the diagonal grows at most quadratically with the input norm
    \[
        k(\bm{x}, \bm{x}) \le C_1 + C_2\mathbb E_{\bm w}[\langle \bm w, \bm x\rangle ^2].
    \]
    Because the weights in both kernels share the same covariance matrix $\frac{1}{d}\bm \Gamma$, we have that
    \[
        \mathbb E_{\bm w}[\langle \bm w, \bm x\rangle ^2] = \frac{1}{d}\bm x^\top \bm \Gamma \bm x\leq \frac{\lambda_{\max}(\bm \Gamma)}{d}\|\bm x\|^2\, ,
    \]
    in both cases. Since $\lambda_{\max}(\bm \Gamma) = A+B$ and $B = \mathcal O\left(\frac{\eta^2}{d}\right)$ we have
    \[
         \mathbb E_{\bm w}[\langle \bm w, \bm x\rangle ^2]\leq C'\frac{\eta^2\|\bm x\|^2}{d^2}
    \]
    for some absolute constant $C' > 0$. Hence, we can find $C_2' > 0$ such that
    \[
        k(\bm{x}, \bm{x}) \le C_1 + C_2'\frac{\eta^2\|\bm x\|^2}{d^2}.
    \]
    By the Cauchy-Schwarz inequality, absorbing all constants into $C > 0$, the off-diagonal terms are bounded by
    \[
        |k(\bm x, \bm x')| \le \sqrt{k(\bm x, \bm x) k(\bm x', \bm x')} \le \sqrt{\left(1 + C\frac{\eta^2\|\bm x\|^2}{d^2}\right)\left(1 + C\frac{\eta^2\|\bm x'\|^2}{d^2}\right)}
    \]
    and using that $\sqrt {ab} \leq (a+b)/2$, absorbing necessary constants into $C$ again, we have
    \[
        |k(\bm x, \bm x')| \leq C \left(1 + \frac{\eta^2\|\bm x\|^2}{d^2} + \frac{\eta^2\|\bm x'\|^2}{d^2}\right)
    \]
    and using the algebraic identity $(a-b)^2 \le 2a^2 + 2b^2$, we get
    \[
        |k_1^*(\bm x, \bm x') - k_1(\bm x, \bm x')|^2 \leq 2|k_1^*(\bm x, \bm x')|^2 + 2|k_1(\bm x, \bm x')|^2 
    \]
    and applying our bound gives
    \[
        |k_1^*(\bm x, \bm x') - k_1(\bm x, \bm x')|^2 \leq 4C \left(1 + \frac{\eta^2\|\bm x\|^2}{d^2} + \frac{\eta^2\|\bm x'\|^2}{d^2}\right)^2\, .
    \]
    Finally, using the identity $(a+b+c)^2 \leq 3(a^2 + b^2 + c^2)$, we have
    \[
        |k_1^*(\bm x, \bm x') - k_1(\bm x, \bm x')|^2 \leq M \left(1 + \frac{\eta^4\|\bm x\|^4}{d^4} + \frac{\eta^4\|\bm x'\|^4}{d^4}\right) \, ,
    \]
    where $M$ is a constant depending on $L_{\sigma}$, $\sigma(0)$ and $C$. Therefore, the integral over the unbounded set is bounded by
    \[
        \iint_{(B_{R, c}\times B_{R, c})^c}|k_1^* - k_1|^2 \mathrm d \rho_{\mathcal X}(\bm x)\mathrm d \rho_{\mathcal X}(\bm x') \leq \iint_{(B_{R, c}\times B_{R, c})^c} M \left(1 +  \frac{\eta^4\|\bm x\|^4}{d^4} +  \frac{\eta^4\|\bm x'\|^4}{d^4}\right)  \mathrm d \rho_{\mathcal X}(\bm x)\mathrm d \rho_{\mathcal X}(\bm x') \, .
    \]
    Noting that
    \begin{align*}
        \Big( B_{R,c} \times B_{R,c}\Big)^c &\subset \left\{\Big\{\|\bm x\| > R \sqrt d\Big\} \times \mathbb R^d\right\} \cup \left\{\Big\{\|\bm x\| < c \sqrt d\Big\} \times \mathbb R^d\right\} \\
        & \quad \cup \left\{\mathbb R^d \times \Big\{\|\bm x'\| > R \sqrt d\Big\}\right\} \cup \left\{\mathbb R^d \times \Big\{\|\bm x'\| < c \sqrt d\Big\}\right\}\, ,
    \end{align*}
    since $\rho_{\mathcal X}$ is a probability measure, by symmetry, using the union bound we have that
    \[
        \iint_{(B_{R, c}\times B_{R, c})^c} 1\mathrm d \rho_{\mathcal X}(\bm x)\mathrm d \rho_{\mathcal X}(\bm x') \leq 2\mathbb P\left(\|\bm x\| > R\sqrt {d}\right) + 2\mathbb P\left(\|\bm x\| < c\sqrt {d}\right)\, .
    \]
    If $\bm x \sim \mathcal{N}(0, \bm I_d)$, by standard Gaussian concentration we have the following probability bound
    \[
        \mathbb{P}\Big(\|\bm x\| > R\sqrt{d}\Big) = \mathbb{P}\Big(\|\bm x\| > \sqrt{d} + (R-1)\sqrt{d}\Big) \leq \exp\left(-\frac{(R-1)^2 d}{2}\right)\,.
    \]
    Furthermore, the probability mass of the excluded inner ball is bounded by
    \[
        \mathbb{P}(\|\bm x\| < c\sqrt{d}) = \mathbb{P}(\|\bm x\|^2 < c^2 d) \le \left[c^2 e^{1-c^2}\right]^{d/2}\, .
    \]
    For every fixed $0<c<1$, we have
    \[
        c^2 e^{1-c^2}<1\, .
    \]
    Defining $y := c^2 e^{1-c^2}$, the bound becomes
    \[
        \mathbb P(\|\bm x\|<c\sqrt d) \leq  y^{d/2}.\, 
    \]
    Since $y <1$, we have $\ln y<0$, and therefore
    \[
        y^{d/2} = \exp\!\left(\frac {d}{2} \ln y\right) = \exp(-c'd)\, ,
    \]
    where
    \[
        c' \coloneqq -\frac12\ln y = \frac12\bigl(c^2-1-2\ln c \bigr) >0\, .
    \]

    Give this, it suffices to estimate the term
    \[
        \iint_{\left\{\{\|\bm x\| > R \sqrt d\} \times \mathbb R^d\right\}}\frac{\eta^4}{d^4}\|\bm x\|^4\mathrm d\rho_{\mathcal X}(\bm x) \mathrm d\rho_{\mathcal X}(\bm x')\, ,
    \]
    and by symmetry the other will follow exactly the same.
    Since the integral does not depend on $\bm x'$, we have
    \[
        \iint_{\left\{\{\|\bm x\| > R \sqrt d\} \times \mathbb R^d\right\}}\frac{\eta^4}{d^4}\|\bm x\|^4\mathrm d\rho_{\mathcal X}(\bm x) \mathrm d\rho_{\mathcal X}(\bm x') = \int_{\{\|\bm x\| > R \sqrt d\}}\frac{\eta^4}{d^4}\|\bm x\|^4\mathrm d\rho_{\mathcal X}(\bm x) \, .
    \]
    By Cauchy Schwarz we have that
    \[
        \int_{\{\|\bm x\| > R \sqrt d\}}\frac{\eta^4}{d^4}\|\bm x\|^4\mathrm d\rho_{\mathcal X}(\bm x)  \leq \frac{\eta^4}{d^4} \sqrt{\mathbb P \Big(\|\bm x\| > R\sqrt d\Big)} \sqrt{\mathbb E_{\bm x}[\|\bm x\|^8]}
    \]
    and since $\|\bm{x}\|^2$ follows a $\chi^2_d$ distribution, we have $\mathbb{E}[\|\bm {x}\|^8] = \mathcal O(d^4)$. Substituting this back and using the concentration of the measure again, we get
    \[
        \int_{\{\|\bm x\| > R \sqrt d\}}\frac{\eta^4}{d^4}\|\bm x\|^4\mathrm d\rho_{\mathcal X}(\bm x) = \mathcal O\left[\frac{\eta^4}{d^2}\exp\left(-\frac{(R-1)^2 d}{4}\right)\right] \, .
    \]
    Analogously, we have
    \[
        \int_{\{\|\bm x\| < c \sqrt d\}}\frac{\eta^4}{d^4}\|\bm x\|^4\mathrm d\rho_{\mathcal X}(\bm x) \leq \frac{\eta^4}{d^4}\sqrt{\mathbb P \Big(\|\bm x\| < c\sqrt d\Big)}\sqrt{\mathbb E_{\bm x}[\|\bm x\|^8]} = \mathcal O\left[\frac{\eta^4}{d^2}\exp(-c'd/2)\right] \, .
    \]
    
    Hence, collecting everything, the integral over $\Big( B_{R, c} \times B_{R, c} \Big)^c$ is bounded by
    \[
        \iint_{(B_{R, c}\times B_{R, c})^c}|k_1^* - k_1|^2 \mathrm d \rho_{\mathcal X}(\bm x)\mathrm d \rho_{\mathcal X}(\bm x') =  \mathcal O\left[\frac{\eta^4}{d^2} \exp \left( -\frac{(R-1)^2d}{4}\right) \right] + \mathcal O\left[\frac{\eta^4}{d^2}\exp(-c'd/2)\right] \, .
    \]
    Since $\eta = \Theta(d^\zeta)$, with $\zeta \in [1/2, 1)$, this gives the bound
    \[
        \iint_{(B_{R, c}\times B_{R, c})^c}|k_1^* - k_1|^2 \mathrm d \rho_{\mathcal X}(\bm x)\mathrm d \rho_{\mathcal X}(\bm x') = \mathcal O\left(d^{4\zeta-2}e^{-K d} \right)\, .
    \]
    for the absolute constant $K = \min \left\{\frac{(R-1)^2}{4}, \frac{c'}{2}\right\} > 0$, which does not depend on $d$. Because of the exponential decay, this is strictly smaller than the bound inside $B_{R, c}\times B_{R, c}$.
    \paragraph{Final bound on the norm:}
    From the previous discussions we know
    \[  
        \iint_{B_{R, c}\times B_{R, c}}|k_1^* - k_1|^2 \mathrm d \rho_{\mathcal X}(\bm x)\mathrm d \rho_{\mathcal X}(\bm x') =  \mathcal O\left(\frac{\eta^4\ln^3 d}{d^5}\right)^2
    \]
    and 
    \[
        \iint_{(B_{R, c}\times B_{R, c})^c}|k_1^* - k_1|^2 \mathrm d \rho_{\mathcal X}(\bm x)\mathrm d \rho_{\mathcal X}(\bm x') = \mathcal O\left(d^{4\zeta-2}e^{-K d} \right)\, .
    \]
    Writing $\|k_1^* - k_1\|_{L^2(\rho_{\mathcal X})\times L^2(\rho_{\mathcal X})}$ as the integral of interest and separating into the integration of both sets
    \[
        \iint|k_1^* - k_1|^2 \mathrm d \rho_{\mathcal X}(\bm x)\mathrm d \rho_{\mathcal X}(\bm x') = \iint_{B_{R, c}\times B_{R, c}}|k_1^* - k_1|^2 \mathrm d \rho_{\mathcal X}(\bm x)\mathrm d \rho_{\mathcal X}(\bm x') + \iint_{(B_{R, c}\times B_{R, c})^c}|k_1^* - k_1|^2 \mathrm d \rho_{\mathcal X}(\bm x)\mathrm d \rho_{\mathcal X}(\bm x')
    \]
    we can combine the bounds over both sets to obtain
    \[
        \|k_1^* - k_1\|_{L^2(\rho_{\mathcal X})\times L^2(\rho_{\mathcal X})} = \mathcal O\left(\frac{\eta^4\ln^3 d}{d^5}\right) \,,
    \]
    and since
    \[
        \|T_1^* - T_1 \|_{\text{op}} \leq \|T_1^* - T_1 \|_{\text{HS}} = \|k_1^* - k_1\|_{L^2(\rho_{\mathcal X})\times L^2(\rho_{\mathcal X})}\, ,
    \]
    the same bound shows the operators are close in norm as well.

    \item \paragraph{Proof of the distribution shift identity}
    We consider the integral operators $T_0: L^2(\rho_{\mathcal X}) \to L^2(\rho_{\mathcal X})$ and $T_1: L^2(\rho_{\mathcal X}) \to L^2(\rho_{\mathcal X})$ 
    \[
        (T_0f)(\bm x) = \int_{\mathcal X} k_0(\bm x, \bm x')f(\bm x') \mathrm d\rho_{\mathcal X}(\bm x') \qquad (T_1f)(\bm x) = \int_{\mathcal X} k_1(\bm x, \bm x')f(\bm x') \mathrm d\rho_{\mathcal X}(\bm x'),
    \]
    and because $T_0$ and $T_1$ are trace-class and compact, Mercer's theorem guarantees that both kernels can be diagonalized. For $k_0$ we have
    \[
        k_0(\bm x, \bm x') = \sum_{i=1}^\infty \xi_i \varphi_i(\bm x)\varphi_i(\bm x')
    \]
    where $\{\varphi_i\}_{i=1}^\infty$ is an orthonormal system and $\{\xi_i\}_{i=1}^\infty$ is a family of eigenvalues associated with these basis such that $\xi_1 \geq \xi_2 \geq ... > 0$. For $k_1$ we have a similar result and
    \[
        k_1(\bm x, \bm x') = \sum_{i=1}^\infty \chi_i \psi_i(\bm x)\psi_i(\bm x')\, ,
    \]
    where $\{\psi_i\}_{i=1}^\infty$ is an orthonormal system and $\{\chi_i\}_{i=1}^\infty$ are the associated non-increasing eigenvalues. In particular, if the kernels are universal, we have that $\overline{\text{span}\{\varphi_i: i \geq 1\}} = \overline{\text{span}\{\psi_i: i \geq 1\}} = L_2(\mathcal X, \rho_{\mathcal X})$.
    
    Consider the matrix from \cref{eq:covz} (without the $\frac{1}{d}$ scaling), written in short form
    \[
        \bm \Gamma  = A\bm I_d + B\bm w^* (\bm w^*)^\top
    \]
    with constants $A$ and $B$ from the original definition satisfying $A + B > |B|$. In particular, given $\bm w \sim \mathcal N\left(0, \frac{1}{d}\bm I\right)$ and $\langle \bm w, \bm x\rangle$, we know that for $\bm z \sim \mathcal N\left(0, \frac{1}{d}\bm \Gamma\right)$ we can write
    \[
        \langle \bm z, \bm x\rangle = \langle \bm w, \bm \Gamma^{1/2} \bm x \rangle
    \]
    and the kernels are related by
    \begin{align*}
        k_1(\bm x, \bm x') &= \mathbb E_{\bm z \sim \mathcal N\left(0, \frac{1}{d}\bm \Gamma\right)}[\sigma(\langle \bm z, \bm x\rangle)\sigma(\langle\bm z, \bm x'\rangle)] \\
        &= \mathbb E_{\bm w \sim \mathcal N\left(0, \frac{1}{d}\bm I_d\right)}[\sigma(\langle \bm w, \bm \Gamma^{1/2} \bm x \rangle)\sigma(\langle \bm w, \bm \Gamma^{1/2} \bm x \rangle)]\\
        &= k_0(\bm \Gamma^{1/2}\bm x, \bm \Gamma^{1/2}\bm x').
    \end{align*}
    
    Let us define the measure $\nu$ as the pushforward measure $\nu = (\bm \Gamma^{\frac{1}{2}})_{\#} V$
    \[
        \nu(V) = (\rho_{\mathcal X} \circ \bm \Gamma ^{-1/2})(V)= \rho_{\mathcal X} (\bm \Gamma^{-1/2}V)\, ,
    \]
    for every measurable set $V$ of $(\bm \Gamma ^{1/2} \mathcal X)$. From now on, we denote the input space transformed by the $\bm \Gamma^{1/2}$ matrix as $(\bm \Gamma^{1/2}\mathcal X) := \mathcal Z \subset \mathbb R^d.$
    
    Consider now, the spaces $L_2(\mathcal X, \rho_{\mathcal X})$ and $L_2(\mathcal Z, \nu)$, and $T_0^{\nu}:\mathcal H_0 \to \mathcal H_0$ the integral operator of $k_0$ w.r.t $\nu$
    \[
        (T_0^{\nu} f)(\bm z) = \int_{\mathcal Z } k_0(\bm z, \bm z')f(\bm z') \mathrm d\nu(\bm z'),\ \forall \bm z \in \mathcal Z
    \]
    
    Due to Mercer's decomposition we have that $k_0$ can be diagonalized into a family of eigenvalues $\{\omega_i\}_{i\in N}$ and eigenfunctions $\{\bm e_i\}_{i \in N}$ which are orthonormal in $L_2(\mathcal Z, \nu)$ and 
    \[
        k_0(\bm z, \bm z') = \sum_{i \in \mathbb N} \omega_i \bm e_i(\bm z) \bm e_i(\bm z') \ \text{ w.r.t. $\nu$},
    \]
    where $\{\omega_i\}$ is not necessarily the same family as $\{\xi_i\}$, nor $\{\bm e_i\}$ are the same eigenfunctions as $\{\bm \varphi_i\}$.
    
    We define $\bm h_i(\cdot) = \bm e_i \circ \bm \Gamma^{1/2}(\cdot)$, the integral operator $T_1$ can be written as
    \begin{align*}
        (T_1 f)(\bm x) &= \int_{\mathcal X} k_1(\bm x, \bm x')f(\bm x') \mathrm d\rho_{\mathcal X}(\bm x') \\
        & = \int_\mathcal X k_0(\bm \Gamma^{1/2}\bm x, \bm \Gamma ^{1/2}\bm x') f(\bm x') \mathrm d\rho_{\mathcal X}(\bm x')\\
        & = \int_\mathcal Z k_0(\bm z, \bm z') (f\circ \bm \Gamma^{-1/2}) ( \bm z') \mathrm d \nu(\bm z'),
    \end{align*}
    and if we plug in $f = \bm h_i$ we get
    \begin{align*}
        (T_1 f)(\bm x) &= \int_\mathcal Z k_0(\bm z, \bm z') (\bm h_i\circ \bm \Gamma^{-1/2}) ( \bm z') \mathrm d \nu(\bm z')\\
        &= \int_\mathcal Z k_0(\bm z, \bm z') (\bm e_i \circ \bm \Gamma ^{1/2}\circ \bm \Gamma^{-1/2}) ( \bm z') \mathrm d \nu(\bm z')\\
        &= \int_\mathcal Z k_0(\bm z, \bm z') \bm e_i( \bm z') \mathrm d \nu(\bm z')\\
        & = \omega_i \bm e_i(\bm z) = \omega_i \bm e_i(\bm \Gamma^{1/2}\bm x),
    \end{align*}
    so $\{\bm e_i\circ \bm \Gamma ^{1/2}\}$ are eigenfunctions of $T_1$ with the same eigenvalues from $k_0$ w.r.t. $\nu$. Also, 
    \[
        \int_\mathcal X \bm e_i(\bm \Gamma^{1/2}\bm x) \bm e_j(\bm \Gamma^{1/2}\bm x) \mathrm d\rho_{\mathcal X}(\bm x) = \int_{\mathcal Z}  \bm e_i(\bm z) \bm e_j(z) \mathrm d\nu(\bm z) = \delta_{i,j},
    \]
    which implies $\{\bm e_i\circ \bm \Gamma ^{1/2}\}$ are orthonormal in $L_2(\mathcal X, \rho_{\mathcal X})$. Therefore, we can also diagonalize $k_1$ such that 
    \[
        k_1(\bm x, \bm x') = \sum_{i\in \mathbb N} \omega_i \bm e_i(\bm \Gamma^{1/2}\bm x)\bm e_i(\bm \Gamma^{1/2}\bm x') \ \text{ w.r.t $\rho_{\mathcal X}$}.
    \]    
\end{proof}

\subsection{Proof of \cref{thm:eigenv-decay}}\label{app:thm-eigenv-decay}

\begin{proof}
    Since we work with the Gaussian measures, the domain is naturally unbounded, thus we split the proof in cases.
    \item \paragraph{Difference under bounded and unbounded domain}
    Supposing we consider an unbounded integration space $\mathcal X$, introduce the bounded set 
    \[
        B_R = \{\bm y \in \mathbb R^d: \|\bm y\| < R\},
    \]
    and we want to show that outside this set we able to control the fluctuations of the eigenvalues well enough.
    
    Since the Gaussian measures are regular, for a given $\varepsilon > 0$, we can choose $R = R(\varepsilon) > 0$ such that
    \[
        \rho_{\bm \Gamma}(B_R^c) < \varepsilon \quad \rho_{\bm I_d}(B_R^c) < \varepsilon
    \]
    and using this we write
    \[
        T_1 f(\bm x) = T_1^{R}f(\bm x) + E_1f(\bm x) = \int_{B_R} k_1(\bm x, \bm x') f(\bm x') \bm{1_{B_R}}(\bm x)\mathrm d\rho_{\bm I_d}(\bm x') + \int_{B_R^c} k_1(\bm x, \bm x') f(\bm x')\mathrm d\rho_{\bm I_d}(\bm x')
    \]
    and similarly $T_0 = T_0^R + E_0$, noting that since $T_1,T_0$ are compact and self adjoint, $T_0^R, T_1^R, E_0, E_1$ must also be.
    
     
    For the bounded domain case, we assume the following result holds: for every $k \geq 0$,
    \[
        c\lambda_k(T_0^R) \leq \lambda_k(T_1^R) \leq C\lambda_k(T_0^R).
    \]
    
    In fact, since the lower bound of the function $r$ does not depend on the norm, its value is the same for the bounded and unbounded case, thus the lower bound on the eigenvalues comes for free
    \[
        c\lambda_k(T_0) \leq \lambda_k(T_1).
    \]
    
    Furthermore, since $\sigma$ is $L_\sigma$-Lipschitz, we have $|\sigma(z) - \sigma(0)| \leq L|z|$ and
    \[
        k_1(\bm{x}, \bm{x}) \leq 2L^2 \mathbb{E}_{\bm{w}}[(\bm{w}^T \bm{x})^2] + 2\sigma(0)^2.
    \]
    Since for $k_1$, $\bm w \sim \mathcal N(0, \frac{1}{d}\bm \Gamma)$, we have
    \[
         k_1(\bm{x}, \bm{x}) \leq C\left(1+\frac{\|\bm x\|^2}{d}\right).
    \]
    for some absolute constant $C > 0$ depending on $\lambda_{\max}(\bm \Gamma), L$ and $\sigma(0)$.
    \[
        \|E_1\|_\op \leq \Trarg{E_1} \leq  \int_{B_R^c} k_1(\bm x, \bm x)\mathrm d\rho_{\bm I_d}(
        \bm x)\leq  \sqrt{\mathbb E_{\bm x}\left[ C\left(1+\frac{\|\bm x\|^2}{d}\right)^2 \right]}\sqrt{\rho_{\bm I_d}(B_R^c)}\, .
    \]
    First, we have $\sqrt{\rho_{\bm I_d}(B_R^c)} < \sqrt \varepsilon$. Also, we can calculate
    \[
        \mathbb E_{\bm x}\left[ C\left(1+\frac{\|\bm x\|^2}{d}\right)^2 \right] =C\mathbb E_{\bm x}\left[1 + \frac{2\|\bm x\|^2}{d}+ \frac{\|\bm x\|^4}{d^2} \right] \leq 6C
    \] 
    hence
    \[
        \|E_1\|_\op \leq \sqrt {6C \varepsilon} \coloneqq C_1' \sqrt{\varepsilon}\,. 
    \]
    Similarly, we have that $\|E_1\|_\op \leq C_0'\sqrt{\varepsilon}$ for some absolute constant $C_0'$.

    Using Weyl's monotonicity theorem we get, for all $k \geq 0$,
    \[
        |\lambda_k(T_1) - \lambda_k(T_1^R)| \leq \|E_1\|_\op < C_1'\sqrt{\varepsilon}
    \]
    and 
    \[
        |\lambda_k(T_0) - \lambda_k(T_0^R)| \leq \|E_0\|_\op < C_0'\sqrt{\varepsilon}
    \]
    which shows the real eigenvalues and the truncated ones are close whenever you fix a radius $R$.     
    
    \item \paragraph{Result for bounded domain}
    
    Based on the above discussion we set out to prove that the eigenvalue decay is maintained in the bounded domain. 
    
    We assume that the integration space $\mathcal X$ is bounded, i.e. $\|\bm x\| < \infty$ for all $\bm x \in \mathcal X$. 
    
    Consider the operator $\mathcal T_\#:  L^2(\rho_{\bm I_d}) \to L^2(\rho_{\bm \Gamma})$ acting on $f$ by the following rule
    \[
        \mathcal T_\# f = f \circ \bm \Gamma^{-1/2}
    \]
    which is designed to push $f$ forward to the transformed space where the new kernel is acting. Since all measures in play are Gaussian, and because $\bm \Gamma$ is full rank, we can perform a change of variables $\bm y = \bm \Gamma ^{1/2} \bm x$ to get
    \[
        \|\mathcal T_\# f\|^2_{L^2(\rho_{\bm \Gamma})} = \int |f\circ \bm \Gamma^{-1/2}(\bm y)|^2 \mathrm d\rho_{\bm \Gamma}(\bm y) = \int |f(\bm x)|^2 \mathrm d\rho_{\bm I_d}(\bm x) = \|f\|^2_{L^2(\rho_{\bm I_d})},
    \]
    thus $\mathcal T_\#$ is an isometry between both spaces and in particular an isomorphism. Similarly, we introduce its conjugate operator $(\mathcal T_\#)^* = \mathcal T_\#^{-1} := \mathcal{T^\#} : L^2(\rho_{\bm \Gamma}) \to L^2(\rho_{\bm I_d})$ which corresponds to the pullback operator
    \[
        \mathcal T^\# g = g\circ \bm \Gamma^{1/2},
    \]
    and is also an isometry between the spaces.
    
    Courant-Fischer's theorem characterizes the $k$-th eigenvalue of any operator by the identity
    \[
        \lambda_k(T_1) = \sup_{V:\dim(V) = k}\inf_{f\in V}\frac{\langle f, T_1 f\rangle}{\|f\|^2_{L^2(\rho_{\bm I_d})}}
    \]
    and we can expand quadratic form into 
    \begin{align*}
        \langle f, T_1 f\rangle_{L^2(\rho_{\bm I_d})} &= \iint f(\bm x)k_1(\bm x, \bm x')f(\bm x')\mathrm d\rho_{\bm I_d}(\bm x)\mathrm d\rho_{\bm I_d}(\bm x')\\
        &=\iint f(\bm \Gamma^{-1/2}\bm y)k_0(\bm y, \bm y')f(\bm \Gamma^{-1/2}\bm y')\mathrm d\rho_{\bm \Gamma}(\bm y)\mathrm d\rho_{\bm \Gamma}(\bm y')\\
        &=\iint f(\bm \Gamma^{-1/2}\bm y)k_0(\bm y, \bm y')f(\bm \Gamma^{-1/2}\bm y')\left[\frac{\mathrm d\rho_{\bm \Gamma}}{\mathrm d\rho_{\bm I_d}}(\bm y)\right]\mathrm d\rho_{\bm I_d}(\bm y)\left[\frac{\mathrm d\rho_{\bm \Gamma}}{\mathrm d\rho_{\bm I_d}}(\bm y')\right]\mathrm d\rho_{\bm I_d}(\bm y')\\
        &:=\iint r(\bm y)f(\bm \Gamma^{-1/2}\bm y)k_0(\bm y, \bm y')r(\bm y')f(\bm \Gamma^{-1/2}\bm y')\mathrm d\rho_{\bm I_d}(\bm y)\mathrm d\rho_{\bm I_d}(\bm y')\\
        &= \langle r\mathcal T_\# f, T_0 (r\mathcal T_\# f)\rangle_{L^2(\rho_{\bm I_d})}
    \end{align*}
    where we use the equality $k_1(\bm x, \bm x') = k_0(\bm \Gamma^{1/2}\bm x, \bm \Gamma^{1/2}\bm x')$ and $r$ is the Radon-Nikodym derivative induced by the two Gaussian measures:
    \[
        r(\bm y) := \frac{\mathrm d\rho_{\bm \Gamma}}{\mathrm d\rho_{\bm I_d}}(\bm y)= \frac{1}{\sqrt{\det \bm \Gamma}}\exp\left(-\frac{1}{2}\bm y^\top (\bm \Gamma^{-1} - \bm I_d)\bm y\right).
    \]
    Given all these definitions, we introduce the operator $\tilde T_1 := \mathcal T_\# T_1 \mathcal T^\#$, leading to the identity
    \[
        \tilde T_1 h(\bm y) = \int k_0(\bm y, \bm y') h(\bm y') \mathrm d\rho_{\bm \Gamma}(\bm y'),
    \]
    furthermore since $T_1$ and $\tilde T_1$ are related by isometries, they are unitarily equivalent and share the same eigenvalues.
    
    To continue we define the operator $M: L^2(\rho_{\bm \Gamma}) \to L^2(\rho_{\bm I_d})$ that maps $g \to r g$ and note that it is bounded and invertible since $r$ is bounded below and above under the assumption that the integration space is bounded, i.e. it is an isomorphism.
    
    If we let $h = M(g)= r g$ we can see
    \begin{align*}
        \langle g, \tilde T_1 g\rangle_{L^2(\rho_{\bm \Gamma})} &= \iint k_0(\bm y, \bm y')g(\bm y)g(\bm y') \mathrm d\rho_{\bm \Gamma}(\bm y)\mathrm d\rho_{\bm \Gamma}(\bm y')\\
        &= \iint k_0(\bm y, \bm y')[r(\bm y)g(\bm y)][r(\bm y')g(\bm y')] \mathrm d\rho_{\bm I_d}(\bm y)\mathrm d\rho_{\bm I_d}(\bm y')
    \end{align*}
    where we use that $r(\bm y)d\rho_{\bm I_d}(\bm y) = d\rho_{\bm \Gamma}(\bm y)$, and thus
    
    \begin{equation}\label{eq:quad_forms}
        \langle g, \tilde T_1 g\rangle_{L^2(\rho_{\bm \Gamma})}= \langle M(g), T_0 M(g)\rangle_{L^2(\rho_{\bm I_d})} = \langle h, T_0 h\rangle_{L^2(\rho_{\bm I_d})}.
    \end{equation}
    
    Analyzing the norms, and writing $g = \frac{1}{r}h$, we get
    \begin{align*}
        \|f\|^2_{L^2(\rho_{\bm I_d})} = \int |f(\bm x)|^2 \mathrm d\rho_{\bm I_d}(\bm x) &= \int |g(\bm y)|^2 \mathrm d\rho_{\bm \Gamma}(\bm y) \\
        &= \int \left|\frac{1}{r(\bm y)}h(\bm y)\right|^2 \mathrm d\rho_{\bm \Gamma}(\bm y)\\
        &= \int \left|\frac{1}{r(\bm x)}h(\bm y)\right|^2 r(\bm y)\mathrm d\rho_{\bm I_d}(\bm y) \\
        &=  \int \frac{1}{r(\bm y)}|h(\bm y)|^2 \mathrm d\rho_{\bm I_d}(\bm y),
    \end{align*}
    and we have that
    \begin{equation}\label{eq:norms}    
        \|f\|^2_{L^2(\rho_{\bm I_d})} = \left\|\frac{1}{\sqrt r}h\right\|^2_{L^2(\rho_{\bm I_d})}.
    \end{equation}
    Thus we can rewrite the variational form as the following identity
    
    \begin{equation}\label{eq:var_forms}
        \frac{\langle f, T_1 f\rangle_{L^2(\rho_{\bm I_d})}}{\|f\|_{L^2(\rho_{\bm I_d})}} = \frac{\langle g, \tilde T_1 g\rangle_{L^2(\rho_{\bm \Gamma})}}{\|g\|^2_{L^2(\rho_{\bm \Gamma})}} = \frac{\langle h, T_0 h\rangle_{L^2(\rho_{\bm I_d})}}{\left\|\frac{1}{\sqrt r} h\right\|^2_{L^2(\rho_{\bm I_d})}}.
    \end{equation}
    
    Since $\bm \Gamma - \bm I_d \succ 0$, we have that $\bm y^\top (\bm \Gamma^{-1} - \bm I_d)\bm y \leq 0$ for all $\bm y \in \mathbb R^d$ and
    \[
        r(\bm y) =  \frac{1}{\sqrt{\det \bm \Gamma}}\exp\left(-\frac{1}{2}\bm y^\top (\bm \Gamma^{-1} - \bm I_d)\bm y\right) \geq \frac{1}{\sqrt{\det \bm \Gamma}}, \ \forall \bm y \in \mathbb R^d.
    \]
    
    And because we assumed $\mathcal X$  to be bounded, we can find constants $c, C > 0$ such that $c\leq r(\bm y) \leq C$ for all $\bm y \in \mathcal X$, thus the relationship between the norms of $f$ and $h$ in \cref{eq:norms} gives
    \[
        \frac{1}{C}\|h\|^2_{L^2(\rho_{\bm I_d)}} \leq \|f\|^2_{L^2(\rho_{\bm I_d)}} = \|g\|^2_{L^2(\rho_{\bm \Gamma)}}\leq \frac{1}{c}\|h\|^2_{L^2(\rho_{\bm Id)}}.
    \]
    
    Combining these bounds with the identity for the quadratic forms (\ref{eq:quad_forms}), we know that for a function $g$ from a fixed subspace $V'$ with dimension $k$, if $h = M(g)$, we have
    \[
        c\frac{\langle h, T_0 h \rangle}{\|h\|^2_{ L^2(\rho_{\bm I_d})}} \leq \frac{\langle g, \tilde T_1 g \rangle}{\|g\|^2_{ L^2(\rho_{\bm \Gamma})}} \leq C\frac{\langle h, T_0 h \rangle}{\|h\|^2_{ L^2(\rho_{\bm I_d})}}.
    \]
    
    Now, given a subspace $V' \subset L^2(\rho_{\bm \Gamma})$ of dimension $k$, we define
    \[
        W = M(V') := \{rg : g \in V'\} \subset L^2(\rho_{\bm I_d}).
    \]
    Because $M$ is bijective and isomorphic, $W$ has dimension $k$ and the mapping $V' \mapsto W$ is a bijection, allowing us to associate every subspace $V' \subset L^2(\rho_{\bm \Gamma})$ with a respective unique subspace $W \subset  L^2(\rho_{\bm I_d})$. 
    
    Thus, taking the infimum over all functions inside $V'$ is the same as taking the infimum over the associated subspace $W$ and 
    \[
        c\left(\inf_{h \in W}\frac{\langle h, T_0 h \rangle}{\|h\|^2_{ L^2(\rho_{\bm I_d})}}\right) \leq \inf_{g \in V'}\frac{\langle g, \tilde T_1 g \rangle}{\|g\|^2_{ L^2(\rho_{\bm \Gamma})}}  \leq C\left(\inf_{h \in W}\frac{\langle h, T_0 h \rangle}{\|h\|^2_{ L^2(\rho_{\bm I_d})}}\right),
    \]
    and finally taking the supremum over all possible sets $V'$ of dimension $k$ is the same as taking the supremum over the associated sets $W$ and we get
    \[
        c\left(\sup_{\dim W = k}\inf_{h \in W}\frac{\langle h, T_0 h \rangle}{\|h\|^2_{ L^2(\rho_{\bm I_d})}}\right) \leq \sup_{\dim V' = k}\inf_{g \in V'}\frac{\langle g, \tilde T_1 g \rangle}{\|g\|^2_{ L^2(\rho_{\bm \Gamma})}}  \leq C\left(\sup_{\dim W = k}\inf_{h \in W}\frac{\langle h, T_0 h \rangle}{\|h\|^2_{ L^2(\rho_{\bm I_d})}}\right),
    \]
    which are exactly the expressions for the eigenvalues of $\lambda(T_0)$ and $\lambda_k(\tilde T_1) = \lambda_k(T_1)$, giving the desired result
    \[
        c \lambda_k(T_0) \leq \lambda_k(\tilde T_1) = \lambda_k(T_1) \leq C\lambda_k(T_0).
    \]

\end{proof}
\subsection{Proof of Theorem \ref{thm:cov-expansion}}\label{app:thm-cov-expansion}
\begin{proof}
    First, we denote $\bm \Lambda := \bm I_d + \frac{\gamma_2}{\gamma_1}\bm u \bm u^\top$. Then by Sherman–Morrison, for $\bm \Sigma = \gamma_1 \bm \Lambda \bm I_d$, we have:
    \[
        \bm \Sigma^{-1} = \frac{1}{\gamma_1}\bm I_d - \frac{\gamma_2}{\gamma_1(\gamma_1+\gamma_2)}\bm u \bm u^\top,
    \]
    and, in particular, the determinants are given by 
    \[
        \det(\bm \Sigma) = \det (\bm \Lambda) \det (\gamma_1 \bm I_d) = \frac{\gamma_1+\gamma_2}{\gamma_1}\det (\gamma_1 \bm I_d).
    \]
    Therefore, we write the expectation of $G$ using the explicit Gaussian density under covariance $\bm \Sigma$ and use the Taylor expansion of the exponential function to obtain
    \begin{align*}
        \mathbb E_{\bm w \sim \mathcal N(0, \bm \Sigma)}[G(\bm w)] &= \int G(\bm w)\frac{1}{(2\pi)^{d/2}\sqrt{\det{\bm \Sigma}}}\exp\left({-\frac{\bm w^\top \bm \Sigma^{-1}\bm w}{2}}\right)\mathrm d\bm w\\
        &=\int \sqrt{\frac{\gamma_1}{\gamma_1+\gamma_2}}\exp\left(\frac{\gamma_2}{2\gamma_1(\gamma_1+\gamma_2)}\langle \bm u, \bm w\rangle^2 \right)\frac{G(\bm w)}{(2\pi)^{d/2}\sqrt{\det \gamma_1\bm I_d}}e^{-\frac{\|\bm w\|^2}{2\gamma_1}}\mathrm d\bm w\\
        &=\sqrt{\frac{\gamma_1}{\gamma_1+\gamma_2}} \int \sum_{k=0}^\infty \frac{1}{k!}\left(\frac{\gamma_2}{2\gamma_1(\gamma_1+\gamma_2)}\right)^k \langle \bm u, \bm w\rangle^{2k} G(\bm w)\mathrm d\rho_{\gamma_1 \bm I_d}(\bm w)\\
        &=\sqrt{\frac{\gamma_1}{\gamma_1+\gamma_2}}\sum_{k=0}^\infty \frac{1}{k!}\left(\frac{\gamma_2}{2\gamma_1(\gamma_1+\gamma_2)}\right)^k\int \langle \bm u, \bm w\rangle^{2k} G(\bm w)\mathrm d\rho_{\gamma_1 \bm I_d}(\bm w)\\
        &= \sqrt{\frac{\gamma_1}{\gamma_1+\gamma_2}}\sum_{k=0}^\infty \frac{1}{k!}\left(\frac{\gamma_2}{2\gamma_1(\gamma_1+\gamma_2)}\right)^k \mathbb E_{\bm w \sim \mathcal N(0, \gamma_1\bm I_d)}[\langle \bm u, \bm w\rangle^{2k}G(\bm w)],
    \end{align*}
    where we used the relationship between the determinants in the second passage.
    
    Next, we justify the application of Fubini's theorem. If we denote $\phi_{\gamma_1}(\bm w) = e^{-\frac{\|\bm w\|^2}{2\gamma_1}}$ and define
    \[
        f(k, \bm w) = \frac{1}{k!}\left(\frac{\gamma_2}{2\gamma_1(\gamma_1+\gamma_2)}\right)^k \langle \bm u, \bm w\rangle^{2k} G(\bm w)\mathrm \phi_{\gamma_1}(\bm w),
    \]
    and we check its sufficient condition: absolute integrability
    \[
        \int \sum_{k \geq 0} |f(k, \bm w)| \mathrm d \bm w < \infty.
    \]

    Looking at this sum we can see
    \begin{align*}
        \sum_{k \geq 0} |f(k, \bm w)| \mathrm d \bm w  &= \sum_{k\geq 0}\frac{1}{k!}\left|\left(\frac{\gamma_2}{2\gamma_1(\gamma_1+\gamma_2)}\right)^k \langle \bm u, \bm w\rangle^{2k} G(\bm w)\mathrm \phi_{\gamma_1}(\bm w)\right| \\
        &\leq \sum_{k\geq 0} \frac{1}{k!}\left|\left(\frac{\gamma_2}{2\gamma_1(\gamma_1+\gamma_2)}\right) \langle \bm u, \bm w\rangle^{2}\right|^k |G(\bm w)||\mathrm \phi_{\gamma_1}(\bm w)| \\
        &= \exp\left(\left|\frac{\gamma_2}{2\gamma_1(\gamma_1+\gamma_2)}\right| \langle \bm u, \bm w\rangle^{2}\right) |G(\bm w)\mathrm| \phi_{\gamma_1}(\bm w)\\
        &= \exp\left(\left|\frac{\gamma_2}{2\gamma_1(\gamma_1+\gamma_2)}\right| \langle \bm u, \bm w\rangle^{2} - \frac{\|\bm w\|^2}{2\gamma_1}\right) |G(\bm w)|.
    \end{align*}

    If this exponent is negative, the Gaussian density is able to suppress the polynomial growth of $G$ under the integral. We always have $\langle \bm u, \bm w \rangle \leq \| \bm w \|$ because of the unit vector $\bm u$, which leads to
    \[
     \exp\left(\left|\frac{\gamma_2}{2\gamma_1(\gamma_1+\gamma_2)}\right|  - \frac{1}{2\gamma_1}\right) = \exp\left(\frac{|\gamma_2| - (\gamma _1 + \gamma_2)}{2\gamma_1(\gamma_1+\gamma_2)}\right)\,.
    \]
    Due to the assumption that $\gamma_1 + \gamma_2 > |\gamma_2|$, the exponent is always negative and consequently
    \[
        \int \sum_{k \geq 0} |f(k, \bm w)| \mathrm d \bm w < \infty.
    \]

    We now study the terms
    \[
         \mathbb E_{\bm w \sim \mathcal N(0, \gamma_1\bm I_d)}[\langle \bm u, \bm w\rangle^{2k}G(\bm w)].
    \]
    By using Stein's Lemma exhaustively we can obtain
    \[
        \mathbb E_{\bm w \sim \mathcal N(0, \gamma_1\bm I_d)}[\langle \bm u, \bm w\rangle^{2k}G(\bm w)] = \sum_{n=0}^k \gamma_1^{k+n}\binom{2k}{2n}(2k -2n-1)!! \mathbb{E}[ D_{\bm u}^{(2n)} G(\bm{w})],
    \]
    and the original expression is given by
    \[
        \mathbb E_{\bm w \sim \mathcal N(0, \bm \Sigma)}[G(\bm w)] = \sqrt{\frac{\gamma_1}{\gamma_1+\gamma_2}}\sum_{k=0}^\infty\frac{1}{k!}\left(\frac{\gamma_2}{2\gamma_1(\gamma_1+\gamma_2)}\right)^k\left[\sum_{n=0}^k \gamma_1^{k+n}\binom{2k}{2n}(2k -2n-1)!! \mathbb{E}[ D_{\bm u}^{(2n)} G(\bm{w})]\right].
    \]
    Since the exchange of the integral and the outer series was justified previously, and for each fixed $k$ the sum over $n$ is finite, the resulting double series is absolutely summable under the same domination bound; hence we may interchange the order of summation.
    \[
        \mathbb E_{\bm w \sim \mathcal N(0, \bm \Sigma)}[G(\bm w)] = \sqrt{\frac{\gamma_1}{\gamma_1+\gamma_2}}\sum_{n=0}^\infty \mathbb{E}[ D_{\bm u}^{(2n)} G(\bm{w})]\left[\sum_{k=n}^\infty\frac{\gamma_1^{n}}{k!}\left(\frac{\gamma_2}{2(\gamma_1+\gamma_2)}\right)^k\binom{2k}{2n}(2k -2n-1)!! \right].
    \]
    Now, due to Lemma \ref{lemma:ith-inf-series} we have the closed form expression for each infinite sum given a fixed $i$, and if we choose $y = \frac{\gamma_2}{2(\gamma_1+\gamma_2)}$:
    \begin{align*}    
        \sum_{k=n}^\infty\frac{\gamma_1^{n}}{k!}\left(\frac{\gamma_2}{2\gamma_1(\gamma_1+\gamma_2)}\right)^k\binom{2k}{2n}(2k -2n-1)!! &= \frac{\gamma_1^{n}}{n!}\left(\frac{\gamma_2}{2(\gamma_1+\gamma_2)}\right)^n\left(\frac{\gamma_1+\gamma_2}{\gamma_1}\right)^{(n+1/2)}\\
        &=\frac{1}{n!}\left(\frac{\gamma_2}{2}\right)^n\sqrt{\frac{\gamma_1+\gamma_2}{\gamma_1}}.
    \end{align*}
    
    Thus, substituting this back into the expansion gives
    \[
        \mathbb E_{\bm w \sim \mathcal N(0, \bm \Sigma)}[G(\bm w)] = \sum_{i=0}^\infty \frac{1}{n!}\left(\frac{\gamma_2}{2}\right)^n\mathbb{E}_{\bm w \sim \mathcal N(0, \gamma_1\bm I_d)}[ D_{\bm u}^{(2n)} G(\bm{w})]
    \]
    and the proof is complete.
\end{proof}
\subsection{Proof of Lemma \ref{lemma:T_R-op-bound}}\label{app:lemma-T_R-op-bound}
\begin{proof}
    
    For $t \in [0, 1]$, we define the parameterized family of covariance matrices
    \[
        \bm \Gamma(t) = A \bm I_d + t B \bm w^* (\bm w^*)^\top,
    \]
    and, if we let $\phi$ be the Gaussian density function induced by $\mathcal N(0, \frac{1}{d} \bm \Gamma(t))$, we define the scalar function $F:[0,1] \to \mathbb R$ given by
    \[
        F(t) = \int G(\bm w)\phi(\bm w, t)\mathrm d \bm w.
    \]
    
    Since $\phi \in C^\infty$, for every $n \geq 0$, the $n$-th derivative of $F$ is well-defined through the distributional property
    \[
        F^{(n)}(t) = \int G(\bm w)\frac{\mathrm d^n}{\mathrm d t^n}[\phi(\bm w, t)]\mathrm d \bm w,
    \]
    and $F \in C^\infty$. Expanding the Taylor series of $F$ around $t=0$ and evaluating at $1$, Taylor's Theorem guarantees the exact identity
    \[
        F(1) = F(0) + F'(0) + \frac{F''(\xi)}{2!}
    \]
    with $\xi \in [0, 1]$.
        
    By Price's Theorem \citep{Price1958, McMahonPrice1964}, we have the formal identity
    \[
        F^{(n)}(t) = \left(\frac{B}{2d}\right)^n\mathbb E_{\bm w \sim \mathcal N(0, \frac{1}{d}\bm \Gamma(t))}[D_{\bm w^*}^{(2n)} G(\bm w) ],
    \]
    therefore
    \[
        F''(\xi) = \left(\frac{B}{2d}\right)^2 \mathbb E_{\bm w \sim \mathcal N(0, \frac{1}{d}\bm \Gamma(\xi))}[D_{\bm w^*}^{(4)} G(\bm w) ],
    \]
    and if we show the fourth derivative under the expectation is bounded independently of the dimension, the multiplying factor will ensure the asymptotic profile of the result.

    Define the quantities
    \[
        h_1 = \langle\bm w, \bm x\rangle \quad h_2 = \langle\bm w, \bm x'\rangle
    \]
    which are joint Gaussian variables such that $(h_1, h_2) := \bm h \sim \mathcal{N}(0, \bm Q_\xi)$ where
    \[
        \bm Q_\xi = \frac{1}{d}\begin{bmatrix}
            \bm x^\top \bm \Gamma(\xi) \bm x & \bm x^\top \bm \Gamma(\xi) \bm x'\\
            \bm x^\top \bm \Gamma(\xi) \bm x' & (\bm x')^\top \bm \Gamma(\xi) \bm x'
        \end{bmatrix}.
    \]
    
    If $\phi_{\bm \Gamma(\xi)}(\bm w)$ denotes the Gaussian density from the measure $\mathcal N(0, \bm \Gamma(\xi))$, let $\phi(h_1, h_2)$ denote the probability density function of this 2D Gaussian.
    
    Then we can write
    \begin{align*}
        \mathbb E_{\bm w \sim \mathcal N(0, \frac{1}{d}\bm \Gamma(t))}[D_{\bm w^*}^{(4)} G(\bm w) ] &= \int D_{\bm w^*}^{(4)} G(\bm w)\phi_{\bm \Gamma(\xi)}(\bm w)\mathrm d \bm w \\
        &= \int D_{\bm w^*}^{(4)} [\sigma(\langle \bm w, \bm x\rangle)\sigma(\langle \bm w, \bm x'\rangle)] \phi_{\bm \Gamma(\xi)}(\bm w)\mathrm d \bm w \\
        &= \iint D_{\bm w^*}^{(4)} [\sigma(h_1)\sigma(h_2)] \phi(h_1, h_2) \mathrm d h_1 \mathrm dh_2.   
    \end{align*}
    
    Next, by the chain rule, applying the directional derivative $D_{\bm w^*} = \langle \bm w^* , \nabla_{\bm w}\rangle$ yields the two dimensional operator over $h_1$ and $h_2$
    \[
        D_{\bm w^*} = \left( \langle \bm w^*, \bm x\rangle  \frac{\partial}{\partial h_1} + \langle \bm w^*, \bm x'\rangle \frac{\partial}{\partial h_2} \right) := \langle \bm U, \nabla_{\bm h}\rangle,
    \]
    where we define
    \[
        \bm U = [\langle \bm w^*, \bm x\rangle, \langle \bm w^*, \bm x'\rangle]^{\top} \in \mathbb R^2.
    \]

    Since $\sigma$ is only Lipschitz, we have no information about its higher order derivatives. Remember that, for any locally integrable $f$ and any smooth compactly supported test function $\varphi$,
    \[
        \langle D_{w^*} f, \varphi \rangle = - \langle f, D_{w^*} \varphi \rangle .
    \]
    Iterating,
    \[
        \langle D_{w^*}^n f, \varphi \rangle = (-1)^n \langle f, D_{w^*}^n \varphi \rangle .
    \]
    Hence, we compute the expectation as a 2D integral and, integrating by parts, we transfer the derivative operator to the density function. 
    \[
        \iint \left[ D_{\bm w^*}^4 \Big( \sigma(h_1)\sigma(h_2) \Big) \right] \phi(h_1, h_2) \, \mathrm d h_1 \mathrm d h_2 =  \iint \sigma(h_1)\sigma(h_2) \left[ D_{\bm w^*}^4 \phi(h_1, h_2) \right] \, \mathrm d h_1 \mathrm d h_2.
    \]

    Now, the closed form of the density is given by
    \[
        \phi(h_1, h_2) = \phi(\bm h) = \frac{1}{2\pi\sqrt{ \det \bm Q(\xi)}} \exp\left( -\frac{1}{2} \bm h^T \bm Q_\xi^{-1} \bm h \right)
    \]
    and differentiating this with our notation gives
    \[
        \nabla_{\bm h}\phi(h_1, h_2) = \nabla_{\bm h}\phi(\bm h) = \phi(\bm h)\nabla_{\bm h}\left( -\frac{1}{2} \bm h^T Q_\xi^{-1} \bm h \right) = -\phi(\bm h) \bm Q_\xi^{-1}\bm h
    \]
    which implies the directional derivative is given by 
    \[
        D_{\bm w^*}\phi(h_1, h_2)  = -\langle \bm U, \bm Q_\xi^{-1}\bm h\rangle \phi(\bm h).
    \]

    Therefore, iterating through this 4 times we obtain 
    \[
        D_{\bm w^*}^4 \phi(h_1, h_2) = \left(\langle \bm U, \bm Q_\xi^{-1} \bm h\rangle^4 - 6\langle \bm U, \bm Q_\xi^{-1} \bm h\rangle^2 \langle \bm U, \bm Q_\xi^{-1} \bm U\rangle + 3\langle \bm U, \bm Q_\xi^{-1} \bm U\rangle^2  \right)\phi(\bm h),
    \]
    and, if we denote the polynomial multiplying $\phi$ by $\mathrm P_4(\bm U, \bm h, \bm Q_{\xi})$, we can see that
    \[
        |\mathrm P_4(\bm U, \bm h, \bm Q_{\xi})| \leq \|\bm U\|^4 \left(\|\bm Q^{-1}_\xi\|^4_{\text{op}}\|\bm h\|^4 + 6 \|\bm Q^{-1}_\xi\|^3_{\text{op}}\|\bm h\|^2 + 3\|\bm Q^{-1}_\xi\|^2_{\text{op}}\right).
    \]

    Using that $\sigma$ is $L_{\sigma}$-Lipschitz we have,
    \[
        |\sigma(t)| \le M_{L_{\sigma}}(1 + |t|)
    \]
    where $M_{L_{\sigma}} = \max(|\sigma(0)|, L_{\sigma})$, therefore the product is bounded by
    \[
        |\sigma(h_1)\sigma(h_2)| \le M_{L_{\sigma}}^2(1 + |h_1|)(1 + |h_2|) \leq M_{L_{\sigma}}^2(2+ \|\bm h\|^2).
    \]

    Thus, the absolute value of $R$ can be bounded with
    \begin{align*}
        |R(\bm x, \bm x')| &\leq \frac{B^2}{8d^2}\iint |\sigma(h_1)\sigma(h_2)| |\mathrm P_4(\bm U, \bm h, \bm Q_{\xi})| \phi(h_1, h_2) \, \mathrm d h_1 \mathrm d h_2\\
        & \leq \frac{B^2}{8d^2}\iint M_{L_{\sigma}}^2(2 + \|\bm h\|^2)  |\mathrm P_4(\bm U, \bm h, \bm Q_{\xi})| \phi(h_1, h_2) \, \mathrm d h_1 \mathrm d h_2\\
        &\leq \frac{B^2}{8d^2}P^*(\bm U, \bm Q_{\xi}),
    \end{align*}
    where we define
    \begin{align*}
        P^*(\bm U, \bm Q_{\xi}) &= M_L^2 \|\bm U\|^4 \cdot \mathbb{E}_{\bm h} \left[ (2 + \|\bm h\|^2)(\|\bm Q_\xi^{-1}\|_{\text{op}}^4\|\bm h\|^4 + 6\|\bm Q_\xi^{-1}\|_{\text{op}}^3\|\bm h\|^2 + 3\|\bm Q_\xi^{-1}\|_{\text{op}}^2) \right]\\
        &=M_L^2 \|\bm U\|^4 \cdot \mathbb{E}_{\bm h} \Big[ \|\bm Q_\xi^{-1}\|_{\text{op}}^4\|\bm h\|^6 + (2\|\bm Q_\xi^{-1}\|_{\text{op}}^4 + 6\|\bm Q_\xi^{-1}\|_{\text{op}}^3)\|\bm h\|^4 \\
        &\qquad +(3\|\bm Q_\xi^{-1}\|_{\text{op}}^2 + 12\|\bm Q_\xi^{-1}\|_{\text{op}}^3)\|\bm h\|^2 + 6\|\bm Q_\xi^{-1}\|_{\text{op}}^2 \Big].
    \end{align*}

    Next, we note that since $\bm h \sim \mathcal N(0, \bm Q_{\xi})$ we know
    \[
        \mathbb{E}[\|\bm{h}\|^{2p}] \leq (2p-1)!! \cdot (2\|\bm Q_\xi\|_{op})^p,
    \]
    therefore distributing the expectation operator and bounding every moment of $\|\bm h\|$ we have
    \begin{align}\label{eq:K'-bound}
        P^*(\bm U, \bm Q_{\xi}) &\leq M_{L_{\sigma}}^2 \|\bm U\|^4 \cdot \Big[ (120 \|\bm Q_\xi^{-1}\|_{\text{op}}^4\|\bm Q_\xi\|_{\text{op}}^3 + 12(2\|\bm Q_\xi^{-1}\|_{\text{op}}^4\| + 6\|\bm Q_\xi^{-1}\|_{\text{op}}^3)\|\bm Q_\xi\|_{\text{op}}^2 \nonumber \\
        &\qquad + 2(2\|\bm Q_\xi^{-1}\|_{\text{op}}^2 + 6\|\bm Q_\xi^{-1}\|_{\text{op}}^3)\|\bm Q_\xi\|_{\text{op}} +4\|\bm Q_\xi^{-1}\|_{\text{op}}^2 \Big].
    \end{align}

    To understand the operator norm of $T_R$ we bound its HS norm. Since the absolute value of $R$ is bounded by $P^*$, we have
    \[
        \|T_R\|^2_{\text{HS}} \leq \frac{B^4}{64d^4}\mathbb E_{\bm x, \bm x'}[P^*(\bm U, \bm Q_{\xi})^2] \, .
    \]
    To bound this value, we first note that, since $\bm w^*$ is a unit vector,
    \[
        \|\bm U\|^{2} = \langle \bm w^*, \bm x\rangle^2 + \langle \bm w^*, \bm x'\rangle^2 := z_1^2 + z_2^2
    \]
    where $z_1, z_2 \sim \mathcal N(0, 1)$ are under i.i.d. Gaussian input distribution and therefore
    \[
        \mathbb E_{\bm x, \bm x'}[\|\bm U\|^k] = \mathbb E[(z_1^2 + z_2^2)^{k/2}] = \mathcal O(1),
    \]
    because the expectation simplifies to a sum of one dimensional Gaussian moments, which are independent of the dimension $d$.
    
    Next, we bound the moments of operator norm $\bm Q_{\xi}$ and its inverse. If we define $\tilde {\bm X} = [\bm x, \bm x'] \in \mathbb R^{d \times 2}$ and let $\tilde {\bm W} = \tilde {\bm X}^\top \tilde {\bm X}$, we have that
    \[
        \bm Q_\xi = \frac{1}{d}\tilde{\bm X}^\top \bm \Gamma(\xi) \tilde{\bm X}
    \]
    and since $\bm \Gamma(\xi)$ is PSD and $A + \xi B > 0$, the eigenvalues of this matrix must be strictly positive. Let $\gamma_{\min} = A$ denote the minimum eigenvalue of $\bm \Gamma(\xi)$. Then, for the inverse matrix we have
    \[
        \bm Q_{\xi}^{-1} \preceq \frac{d}{\gamma_{\min}} \tilde{\bm W}^{-1}.
    \]
    which implies the following bounds on the moments
    \[
        \mathbb{E}_{\bm x, \bm x'} \left[ \|\bm Q_\xi^{-1}\|_{\text{op}}^k \right] \le \left(\frac{d}{\gamma_{\min}}\right)^k \mathbb{E}_{\bm W} \left[ \|\tilde{\bm W}^{-1}\|_{\text{op}}^k \right].
    \]

    Because $\tilde{\bm X}$ is a $d \times 2$ matrix of independent $\mathcal{N}(0, 1)$ entries, the inner product matrix follows a Wishart distribution: $\tilde{\bm W} \sim \mathcal{W}_2(\bm I_2, d)$ and, consequently, its inverse follows an Inverse-Wishart distribution $\tilde{\bm W}^{-1} \sim \mathcal{W}^{-1}_2(\bm I_2, d)$.

    For an $n \times n$ Inverse-Wishart matrix with $d$ degrees of freedom, the $k$-th moments are strictly finite and integrable as long as the degrees of freedom exceed the matrix size, i.e. $d > n + 2k - 1$. In our case, $\tilde{\bm W}$ is $2 \times 2$, therefore we have
    \[
       \mathbb{E}_{\tilde{\bm W}} \left[ \|\tilde{\bm W}^{-1}\|_{\text{op}}^k \right] = \mathcal O\left(\frac{1}{d^k}\right)
    \]
    as long as $d > 2k + 1.$

    As a consequence, since $A$ does not depend on $d$, for a fixed $k$, as the dimension grows, the operator norm is bounded independently of $d$, i.e.
    \[
        \mathbb{E}_{\bm x, \bm x'} \left[ \|\bm Q_\xi^{-1}\|_{\text{op}}^k \right] = \mathcal O(1)\, .
    \]
    To bound the non-inverse moments, we decompose the quadratic form directly
    \[
        \bm{Q}_\xi = \frac{1}{d} \tilde{\bm{X}}^\top (A\bm{I}_d + B\bm{w^*}(\bm{w^*})^\top) \tilde{\bm{X}} = \frac{A}{d} \tilde{\bm{W}} + \frac{B}{d} \bm{z}\bm{z}^\top
    \]
    where $\bm{z} = \tilde{\bm{X}}^\top \bm{w^*} \in \mathbb{R}^2$. Because the columns of $\tilde{\bm{X}}$ are standard Gaussian and $\bm{w^*}$ is a deterministic unit vector, the projection $\bm{z}$ is distributed as a standard 2-dimensional Gaussian, $\bm{z} \sim \mathcal{N}(0, \bm{I}_2)$.

    Applying the triangle inequality to the operator norm we obtain
    \[
        \|\bm{Q}_\xi\|_{\text{op}} \le \frac{A}{d} \|\widetilde{\bm{W}}\|_{\text{op}} + \frac{B}{d} \|\bm{z}\bm{z}^\top\|_{\text{op}} = \frac{A}{d} \|\widetilde{\bm{W}}\|_{\text{op}} + \frac{B}{d} \|\bm{z}\|_2^2\, .
    \]
    Raising the norm to the $k$-th power and applying the inequality $(a+b)^k \le 2^{k-1}(a^k + b^k)$ yields
    \[
        \mathbb{E}_{\bm x, \bm x'}[\|\bm{Q}_\xi\|_{\text{op}}^k] \le 2^{k-1} \left( \frac{A^k}{d^k} \mathbb{E}[\|\widetilde{\bm{W}}\|_{\text{op}}^k] + \left(\frac{B}{d}\right)^k \mathbb{E}[\|\bm{z}\|_2^{2k}] \right)
    \]
    Given $B = \Theta(\eta^2/d)$, with $\eta = \Theta(d^\zeta)$ and $\zeta \in [1/2, 1)$, we have $B = \Theta(d^{2\zeta - 1})$ and the coefficient scales as $\frac{B}{d} = \Theta(d^{2\zeta-2})$. Because $\zeta < 1$ by assumption, it follows that $2\zeta - 2 < 0$ and $\frac{B}{d} = o_d(1)$. Furthermore, because $\|\bm{z}\|_2^2$ follows a chi-squared distribution with 2 degrees of freedom, its moments $\mathbb{E}[\|\bm{z}\|_2^{2k}]$ are constants dependent only on $k$, bounded by $\mathcal{O}(1)$. Hence, we also have
    \[
        \mathbb{E}_{\bm x, \bm x'} \left[ \|\bm Q_\xi\|_{\text{op}}^k \right] = \mathcal O(1)\, .
    \]
    
    Furthermore, for fixed exponents $k_1, k_2 \geq 0$, by Cauchy-Schwarz we are able to obtain
    \[
        \mathbb E_{\bm x, \bm x'}\left[\|\bm Q_{\xi}^{-1}\|_{\text{op}}^{k_1}\|\bm Q_{\xi}\|_{\text{op}}^{k_2}\right] = \mathcal O(1).
    \]    
    
    Lastly, notice that if we square the bound in \ref{eq:K'-bound}, we obtain a combination of powers of $\|\bm U\|, \|\bm Q_{\xi}^{-1}\|_{\text{op}}$ and $\|\bm Q_{\xi}\|_{\text{op}}$. Applying Holder's inequality to decouple the expectations on each of the terms of the inequality and collecting all the bounds we have on the moments of these quantities we get
    \[
        \mathbb E_{\bm x, \bm x'}[P^*(\bm U, \bm Q_\xi)^2] = \mathcal O(1).
    \]
    
    Thus, we have the following result for the HS norm
    \begin{align*}
        \|T_R\|_{\text{HS}}^2 &= \iint |R(\bm x, \bm x')|^2\mathrm d \rho_{\mathcal X}(\bm x)\mathrm d \rho_{\mathcal X}(\bm x')\\
        &\leq  \frac{B^4}{64d^4} \mathbb E_{\bm x, \bm x'}[P^*(\bm U, \bm Q_{\xi})^2]\\
        &=\mathcal O\left(\frac{B^4}{d^4}\right)
    \end{align*}
    which implies
    \[
        \| T_R\|_{\text{op}} \leq \|T_R\|_{\text{HS}} = \mathcal O\left(\frac{B^2}{d^2}\right)
    \]
    and the proof is completed.
\end{proof}
\subsection{Proof of Lemma \ref{lemma:iso-eigenfn}}
\begin{proof}
    We write the ReLU kernel as its closed-form
    \[
        k_0(\bm x, \bm x') = \frac{\|\bm x\|\|\bm x'\|}{2 \pi d}\left[\gamma(\pi -\arccos (\gamma)) +\sqrt{1 - \gamma^2}\right] := \frac{r r'}{2 \pi d}J_1(\gamma)
    \]
    where $r = \|\bm x\|$ and $r' = \|\bm x'\|$ and $J_1$ is the standard first-order arc cosine kernel.
    
    In this context, the function $J_1(\gamma)$ is a dot-product kernel on the unit sphere and, by the Funk-Hencke's theorem, $J_1$ can be expanded into the spherical harmonics basis. Therefore, if we denote $\bm \omega = \frac{\bm x}{\|\bm x\|}$ and $\bm \omega' = \frac{\bm x'}{\|\bm x'\|}$, we have
    \[
        k_0(\bm x, \bm x') = \frac{r r'}{2 \pi d} \sum_{k=0}^{\infty} \lambda_k \sum_{m} Y_{k,m}(\bm \omega) Y_{k,m}(\bm \omega').
    \]
    By moving the $r$ and $r'$ inside the sum, we get the exact Mercer expansion over the Gaussian space
    \[
        k_0(\bm x, \bm x') = \frac{1}{2 \pi}\sum_{k=0}^{\infty} \sum_{m} \lambda_k \left( \frac{r}{\sqrt {d}} Y_{k,m}(\bm \omega) \right) \left( \frac{r'}{\sqrt{d}} Y_{k,m}(\bm \omega') \right).
    \]
    Note that these functions are indeed normalized and pairwise orthogonal since the inner product of two elements of this family can always be separated into an integral acting only on the radii and an integral acting only on the angles. Because the spherical harmonics are pairwise orthogonal, the angular integral will enforce the orthogonality conditions.
    
    Therefore, since this expansion fully constructs the kernel, the only eigenfunctions with non-zero eigenvalues are of the form $\frac{r}{\sqrt d} Y_{k,m}(\bm \omega)$.
\end{proof}
\subsection{Proof of Lemma \ref{lemma:approx-first-order-op}}\label{app:lemma-first-order-op}
\begin{proof}
    For this proof we write
    \begin{align*}
        T_1 f(\bm x) &= AT_0 f(\bm x) + \frac{B}{2d} T_{S} f(\bm x)+  T_R f(\bm x)\\
        &= AT_0 f(\bm x) + \frac{B}{2d} \left(T_{(1*)} f(\bm x) + T_{(2*)} f(\bm x)\right) +  T_R f(\bm x)
    \end{align*}
    and, as we have seen by Lemma \ref{lemma:T_R-op-bound}, $\|T_Rf\|_{L^2(\rho_\mathcal X)} = o_d\left(\frac{|B|}{d}\right)$. Thus, we will study the terms $T_{(1*)}f$ and $T_{(2*)} f$.
    
    {\bf Analysis of the first term:}
    We start by analyzing
    \[
        T_{(1*)}f(\bm x) =\langle \bm x, \bm w^*\rangle\int \frac{(\pi - \theta_{\bm x, \bm x'})}{\pi}\langle \bm x', \bm w^*\rangle f(\bm x')\mathrm d\rho_{\mathcal X}(\bm x'),
    \]    
    and focus on the term under the integral. Note that $f$ and $\langle \cdot, \bm w^*\rangle$ are normalized functions such that $\|f\|_{L^2(\rho_{\mathcal X})} = \|\langle \cdot, \bm w^*\rangle\|_{L^2(\rho_{\mathcal X})} = 1.$
    
    If we consider the set 
    \[
        \mathcal A = \left\{\bm x' \in \mathbb R^d: \left|\theta_{\bm x, \bm x'} - \frac{\pi}{2}\right|< \epsilon \right\}, \epsilon \in \left(0, \frac{\pi}{2}\right),
    \]
    we split the integral over the entire space into the sum of integrals over $\mathcal A$ and $\mathcal A^c$.
    
    
    By Lemma \ref{lemma:integral-concentration-1}, if we let $g = \langle \cdot, \bm w^*\rangle f$, for a given $\epsilon > 0$, the concentration of the Gaussian measure implies
    \[
        \left|\int_{\mathcal A^c} \frac{(\pi - \theta_{\bm x, \bm x'})}{\pi} \langle \bm x', \bm w^*\rangle f(\bm x') \mathrm d\rho_{\mathcal X}(\bm x')\right| \leq \sqrt 2e^{-c_1d\epsilon^2/2}\|g\|_{L^2(\rho_{\mathcal X})}
    \]
    for an absolute constant $c_1 > 0$, provided $g \in L^2(\rho_{\mathcal X})$.
    We can see that
    \[
        \int |\langle \bm x', \bm w^*\rangle f(\bm x')|^2\mathrm d \rho_{\mathcal X}(\bm x')  = \int |\|\bm x'\|^2 (\cos\theta_{\bm x', \bm w^*})^2 f(\bm x')|^2\mathrm d \rho_{\mathcal X}(\bm x') \leq \int \|\bm x'\|^2 |f(\bm x')|^2\mathrm d \rho_{\mathcal X}(\bm x').
    \]
    Now, due to Lemma \ref{lemma:iso-eigenfn}, the orthonormal eigenbasis is of the separated form
    \[
        \psi_{k,m}(\bm{x}) = \frac{\|\bm x\|}{\sqrt d}Y_{k,m}(\bm \omega),
    \]
    a purely radial function governing the magnitude and $Y_{k,m}(\bm \omega)$ is a spherical harmonic of degree $k$, governing the direction $\bm \omega = \frac{\bm x}{\|\bm x\|}$.

    Although the ReLU kernel is not universal on the Gaussian function space, by assumption we can write $f$ as a linear combination of elements of the eigenbasis of the kernel
    \[
        f(\bm x) = \sum_{k,m \geq 0} c_{k, m}\psi_{k, m}(\bm x)
    \]
    therefore
    \begin{align*}
        |f(\bm{x})|^2 &= \left( \sum_{k,m} c_{k,m} \frac{\|\bm{x}\|}{\sqrt{d}} Y_{k,m}(\bm \omega) \right) \left( \sum_{k',m'} \overline{c_{k',m'}} \frac{\|\bm{x}\|}{\sqrt{d}} \overline{Y_{k',m'}(\bm \omega)} \right) \\
        &= \frac{\|\bm{x}\|^2}{d} \sum_{k,m} \sum_{k',m'} c_{k,m} \overline{c_{k',m'}} Y_{k,m}(\bm \omega) \overline{Y_{k',m'}(\bm \omega)}
    \end{align*}

    Decomposing $f$ into this form inside of the integral yields the separation
    \begin{align*}
        \int \|\bm x'\|^2 |f(\bm x')|^2\mathrm d \rho_{\mathcal X}(\bm x') &= \sum_{k,m} \sum_{k', m'} c_{k,m} c_{k',m'} \left(\frac{1}{d}\int_{\mathbb R}| r^2|^2 \mathrm \phi(r)dr\right) \left(\int_{\Omega}Y_{k,m}( \omega)\overline{Y_{k',m'}(\omega)}\mathrm d\Omega\right),\\
        &= \sum_{k,m}|c_{k,m}|^2\left(\frac{1}{d}\int_{\mathbb R}| r^2|^2 \mathrm \phi(r)dr\right)\\
        &=\left(\frac{1}{d}\int_{\mathbb R}| r^2|^2 \mathrm \phi(r)dr\right)
    \end{align*}
    where we denote by $\phi(r)$ the Gaussian density associated with the radial decomposition of the norm and we use Parseval's identity to simplify
    \[
        \sum_{k,m}|c_{k,m}|^2 = \|f\|_{L^2(\rho_{\mathcal X})} = 1.
    \]
    
    The second term of the product is simply the norm of the spherical harmonics, which evaluates to $1$ since it is a normalized function. The first term is the integral of a fourth Gaussian moment so we solve analytically 
    \[
        \frac{1}{d}\int r^4\phi(r)\mathrm dr = \frac{1}{d}\mathbb E[\|\bm x\|^4] = \frac{d(d+2)}{d} = d+2.
    \]
    Therefore, if we define
    \[
        E_1^{(1)}(\bm x) := \langle \bm x, \bm w^*\rangle \int_{\mathcal A^c} \frac{(\pi - \theta_{\bm x, \bm x'})}{\pi}\langle \bm x', \bm w^*\rangle f(\bm x')\mathrm d\rho_{\mathcal X}(\bm x').
    \]
    we have
    \[
        \left|E_1^{(1)}(\bm x)\right| = \mathcal O(de^{-d\epsilon^2})|\langle\bm x, \bm w^*\rangle|
    \]
    and we can see the $L^2$ norm respects
    \[
        \|E_1^{(1)}\|_{L^2(\rho_{\mathcal X})} = \mathcal O(de^{-d\epsilon^2}).
    \]
    
    Moreover, inside of $\mathcal{A}$ we write
    \begin{align*}
        &\langle\bm x, \bm w^*\rangle\int_{\mathcal{A}} \frac{(\pi - \theta_{\bm x, \bm x'})}{\pi}\langle \bm x', \bm w^*\rangle f(\bm x')\mathrm d\rho_{\mathcal X}(\bm x') \\
        & \quad=\frac{\langle\bm x, \bm w^*\rangle}{2} \left[ \int_{\mathcal{A}} \langle \bm x', \bm w^*\rangle f(\bm x')\mathrm d\rho_{\mathcal X}(\bm x') - \int_{\mathcal{A}}  \frac{\theta_{\bm x, \bm x'} - \pi/2}{\pi}\langle \bm x', \bm w^*\rangle f(\bm x')\mathrm d\rho_{\mathcal X}(\bm x')\right]\\
        &\quad := \frac{\langle\bm x, \bm w^*\rangle}{2}\int_{\mathcal{A}} \langle \bm x', \bm w^*\rangle f(\bm x')\mathrm d\rho_{\mathcal X}(\bm x') + E_2^{(1)}(\bm x).
    \end{align*}
    
    Since inside of $\mathcal{A}$ the difference between the angles is bounded by $\epsilon$, using Cauchy-Schwarz we can bound $E_2^{(1)}(\bm x)$ by
    \begin{align*}
        |E_2(\bm x)^{(1)}| &= \left|\langle\bm x, \bm w^*\rangle \int_{\mathcal{A}}  \frac{\theta_{\bm x, \bm x'} - \pi/2}{\pi}\langle \bm x', \bm w^*\rangle f(\bm x')\mathrm d\rho_{\mathcal X}(\bm x')\right| \leq \epsilon|\langle\bm x, \bm w^*\rangle| \int_{\mathcal{A}} |\langle \bm x', \bm w^*\rangle f(\bm x')|\mathrm d\rho_{\mathcal X}(\bm x') \\
        &\leq \epsilon |\langle\bm x, \bm w^*\rangle| \| \langle \cdot, \bm w^*\rangle \|_{L^2(\rho_{\mathcal X})}\| f\|_{L^2(\rho_{\mathcal X})} \leq \epsilon|\langle\bm x, \bm w^*\rangle|,        
    \end{align*}
    because $\langle \cdot, \bm w^*\rangle$ and $f$ are normalized in $L^2(\rho_{\mathcal X})$. Again, this implies the $L^2$ norm respects
    \[
        \|E_2^{(1)}\|_{L^2(\rho_{\mathcal X})} = \mathcal O(\epsilon).
    \]
    
    This gives the following expression for the first term
    \[
        \langle \bm x, \bm w^*\rangle\int \frac{(\pi - \theta_{\bm x, \bm x'})}{\pi}\langle \bm x', \bm w^*\rangle f(\bm x')\mathrm d\rho_{\mathcal X}(\bm x') = \frac{\langle \bm x, \bm w^*\rangle}{2}\int_{\mathcal{A}} \langle \bm x', \bm w^*\rangle f(\bm x')\mathrm d\rho_{\mathcal X}(\bm x') + E_1^{(1)}(\bm x) + E_2^{(1)}(\bm x)
    \]
    
    We can calculate the integral over $\mathcal{A}$ by writing
    \[
        \int_{\mathcal{A}} \langle \bm x', \bm w^*\rangle f(\bm x')\mathrm d\rho_{\mathcal X}(\bm x') = \int \langle \bm x', \bm w^*\rangle f(\bm x')\mathrm d\rho_{\mathcal X}(\bm x') - \int_{\mathcal{A}^c} \langle \bm x', \bm w^*\rangle f(\bm x')\mathrm d\rho_{\mathcal X}(\bm x')
    \]
    and by using the same argument as before, the integral over $\mathcal{A}^c$ induces a function $E_3^{(1)}$ and
    \begin{align*}
        \langle \bm x, \bm w^*\rangle\int_{\mathcal{A}} \langle \bm x', \bm w^*\rangle f(\bm x')\mathrm d\rho_{\mathcal X}(\bm x') &= \langle \bm x, \bm w^*\rangle\int \langle \bm x', \bm w^*\rangle f(\bm x')\mathrm d\rho_{\mathcal X}(\bm x') +  E_3^{(1)}(\bm x)\\
        &=\langle \bm x, \bm w^*\rangle\langle\langle \cdot, \bm w^*\rangle, f\rangle +  E_3^{(1)}(\bm x).
    \end{align*}
    such that $\|E_3^{(1)}\|_{L^2(\rho_{\mathcal X})} = \mathcal O(de^{-d\epsilon^2})$.
    
    Finally, the first term can be written as
    \begin{align*}
        T_{(1*)}f(\bm x) &= \langle \bm x, \bm w^*\rangle\int \frac{(\pi - \theta_{\bm x, \bm x'})}{\pi}\langle \bm x', \bm w^*\rangle f(\bm x')\mathrm d\rho_{\mathcal X}(\bm x') \\
        &= \frac{\langle \bm x, \bm w^*\rangle}{2}\langle\langle \cdot, \bm w^*\rangle,f\rangle + E_1^{(1)}(\bm x) + E_2^{(1)}(\bm x) + E_3^{(1)}(\bm x)\\
        &:= T_{S}^{(1*)}(\bm x)+ E_1^{(1)}(\bm x) + E_2^{(1)}(\bm x) + E_3^{(1)}(\bm x)
    \end{align*}
    and choosing $\epsilon = d^{-1/3}$ we have the bound on the $L^2$ difference by noticing
    \begin{align*}
        \|T_{(1*)}f - T_{S}^{(1*)}f\|_{L^2(\rho_{\mathcal X})} &\leq \left(\|E_1^{(1)}\|_{L^2(\rho_{\mathcal X})} + \|E_2^{(1)}\|_{L^2(\rho_{\mathcal X})} +\|E_3^{(1)}\|_{L^2(\rho_{\mathcal X})}\right) \\
        &\leq \left(\mathcal O(de^{-d^{1/3}})+\mathcal O(d^{-1/3}) + \mathcal O(de^{-d^{1/3}})\right)\\
        &= o_d(1).
    \end{align*} 

    {\bf Analysis of the second term:}
    Now we turn to the second term
    \[
        T_{(2*)}f(\bm x) = \int\frac{\sin \theta_{\bm x, \bm x'}}{2\pi}\left(\frac{\|\bm x'\|}{\|\bm x\|}\langle \bm x, \bm w^*\rangle^2 + \frac{\|\bm x\|}{\|\bm x'\|}\langle \bm x', \bm w^*\rangle^2\right)f(\bm x')\mathrm d\rho_{\mathcal X}(\bm x'),
    \]
    and apply the same treatment.

    Considering the same set $\mathcal A$, we split the integral over the entire space into $\mathcal A$ and $\mathcal A^c$, and first we show the following function is negligible
    \[
        E^{(2)}_1 (\bm x) := \frac{\langle \bm x, \bm w^*\rangle^2}{\|\bm x\|}\int_{\mathcal A^c} \frac{\sin \theta_{\bm x, \bm x'}}{2\pi}\|\bm x'\| f(\bm x')\mathrm d\rho_{\mathcal X}(\bm x') + \|\bm x\|\int_{\mathcal A^c} \frac{\sin \theta_{\bm x, \bm x'}}{2\pi}\frac{\langle \bm x', \bm w^*\rangle^2}{\|\bm x'\|} f(\bm x')\mathrm d\rho_{\mathcal X}(\bm x')
    \]

    Looking at the integral in the second term, again by Lemma  \ref{lemma:integral-concentration-1}, if we define $g = \frac{\langle \cdot, \bm w^*\rangle ^2}{\|\bm \cdot\|}f$ we have that
    \[
        \left|\int_{\mathcal A^c} \frac{\sin \theta_{\bm x, \bm x'}}{2\pi} {\frac{\langle \bm x', \bm w^*\rangle ^2}{\|\bm x'\|} f(\bm x') }\mathrm d \rho_{\mathcal X}(\bm x')\right| \leq \sqrt 2e^{-c_2d\epsilon^2/2}\|g\|_{L^2(\rho_{\mathcal X})},
    \]
    for an absolute constant $c_2 > 0$, as long as $g \in L^2(\rho_{\mathcal X}).$ Notice that, outside of a set with zero measure, we have
    \[
        \langle \bm x', \bm w^*\rangle = \|\bm x'\| \cos \theta_{\bm x', \bm w^*} \implies \frac{\langle \bm x', \bm w^*\rangle ^2}{\|\bm x'\|} = \cos \theta_{\bm x', \bm w^*} \langle \bm x', \bm w^*\rangle
    \]
    therefore calculating the $L^2(\rho_\mathcal X)$ norm of $g$ gives:
    \[
        \int \left|\frac{\langle \bm x', \bm w^*\rangle ^2}{\|\bm x'\|} f(\bm x')\right|^2 \mathrm d \rho_{\mathcal X}(\bm x')= \int \left| \cos \theta_{\bm x, \bm w^*} \langle \bm x', \bm w^*\rangle f(\bm x')\right|^2 \mathrm d \rho_{\mathcal X}(\bm x') \leq \int \left| \langle \bm x', \bm w^*\rangle f(\bm x')\right|^2 \mathrm d \rho_{\mathcal X}(\bm x').
    \]
    Recall that in the analysis of the first term we showed that $\|\cdot \|f \in L^2(\rho_{\mathcal X})$, and since this function dominates $g$, we must have $g \in L^2(\rho_{\mathcal X})$ and the concentration bound holds.
    
    Similarly, if we let $g' = \|\cdot \| f$, we already showed $g' \in L^2(\rho_{\mathcal X})$, and using the lemma again implies
    \[
        \left|\int_{\mathcal A^c} \frac{\sin \theta_{\bm x, \bm x'}}{2\pi} \|\bm x'\|f(\bm x')\mathrm d \rho_{\mathcal X}(\bm x')\right| \leq \sqrt 2e^{-c_2d\epsilon^2/2}\|g'\|_{L^2(\rho_{\mathcal X})}.
    \]
    
    Thus, outside of a set of measure zero, we have the following bound
    \[
        |E_1^{(2)}(\bm x)| \leq \mathcal O(de^{-d \epsilon^2})\left(\|\bm x\| + \frac{\langle \bm x, \bm w^*\rangle ^2}{\|\bm x\|}\right)
    \]
    which implies the $L^2(\rho_{\mathcal X})$ norm is bounded by
    \[
        \|E_1^{(2)}\|_{L^2(\rho_{\mathcal X})} \leq  \mathcal O(de^{-d \epsilon^2})\left(\Big\|\|\cdot\|\Big\|_{L^2} +\left\|\frac{\langle \cdot, \bm w^*\rangle ^2}{\|\cdot\|}\right\|_{L^2} \right) = \mathcal O(d^{3/2} e^{-d\epsilon^2}).
    \]
    
    Next, inside of $\mathcal{A}$,  we have
    \begin{align*}
        &\int_{\mathcal A} \frac{\sin \theta_{\bm x, \bm x'}}{2\pi}\left(\frac{\|\bm x'\|}{\|\bm x\|}\langle \bm x, \bm w^*\rangle^2 + \frac{\|\bm x\|}{\|\bm x'\|}\langle \bm x', \bm w^*\rangle^2\right)f(\bm x')\mathrm d\rho_{\mathcal X}(\bm x') \\
        &\quad := \frac{1}{2 \pi}\int_{\mathcal A}\left(\frac{\|\bm x'\|}{\|\bm x\|}\langle \bm x, \bm w^*\rangle^2 + \frac{\|\bm x\|}{\|\bm x'\|}\langle \bm x', \bm w^*\rangle^2\right)f(\bm x')\mathrm d\rho_{\mathcal X}(\bm x') + E_2^{(2)}(\bm x).
    \end{align*}
    where we define
    \[
        E_2^{(2)}(\bm x) = \int_{\mathcal A} \left[\frac{\sin \theta_{\bm x, \bm x'}-1}{2\pi}\right] \left(\frac{\|\bm x'\|}{\|\bm x\|}\langle \bm x, \bm w^*\rangle^2 + \frac{\|\bm x\|}{\|\bm x'\|}\langle \bm x', \bm w^*\rangle^2\right)f(\bm x')\mathrm d\rho_{\mathcal X}(\bm x').
    \]

    Note that
    \[
        |\sin \theta - 1| = 1- \sin \theta  = 1 - \cos(\theta - \pi/2),
    \]
    and if we expand $\cos x$ using Taylor's theorem around $0$, noting that $\cos'' x = -\cos x$, we obtain
    \[
        \cos x = 1  + \frac{\cos''(c)}{2}x^2 \implies 1 - \cos x = \frac{\cos(c)}{2}x^2
    \]
    for some $c \in [0, 2\pi]$. Since $\cos \theta \leq 1$ for all $\theta \in [0, 2\pi]$ we have the identity
    \[
        1 - \cos x \leq \frac{x^2}{2}.
    \]
    Thus, letting $x = \theta - \pi/2$ we get
    \[
        |\sin \theta - 1| \leq \frac{1}{2}\left(\theta - \frac{\pi}{2}\right)^2, \forall \theta \in [0, 2\pi],
    \]
    and, since the difference between the angles is bounded by $\epsilon$ inside of $\mathcal A$, using Cauchy-Schwarz again we can obtain the following bound for $E_2^{(2)}$:
    \begin{align*}
        |E_2^{(2)}(\bm x)| &\leq \frac{\epsilon^2}{4\pi}\sqrt{\int \left(\frac{\|\bm x'\|}{\|\bm x\|}\langle \bm x, \bm w^*\rangle^2 + \frac{\|\bm x\|}{\|\bm x'\|}\langle \bm x', \bm w^*\rangle^2\right)^2\mathrm d\rho_{\mathcal X}(\bm x')} \sqrt{\int|f(\bm x')|^2\mathrm d\rho_{\mathcal X}(\bm x')}\\
        &\leq \frac{\epsilon^2}{4 \pi}\left(\Big\|\|\cdot\|\Big\|_{L^2(\rho_{\mathcal X})}\frac{\langle \bm x, \bm w^*\rangle^2}{\|\bm x\|} + \|\bm x \|\left\|\frac{\langle \cdot, \bm w^*\rangle^2}{\|\cdot\|}\right\|_{L^2(\rho_{\mathcal X})}\right)
    \end{align*}
    almost everywhere.

    Taking the $L^2(\rho_{\mathcal X})$ norm gives
    \[
        \|E_2^{(2)}(\bm x)\|_{L^2(\rho_{\mathcal X})} \leq \frac{2\epsilon^2}{4\pi} \Big\|\|\cdot\|\Big\|_{L^2(\rho_{\mathcal X})}\left\|\frac{\langle \cdot, \bm w^*\rangle^2}{\|\cdot\|}\right\|_{L^2(\rho_{\mathcal X})} \leq \frac{2\epsilon^2 \sqrt{d}}{4\pi}.
    \]

    Lastly, we are left with
    \[
        T_{(2*)}f(\bm x) = \frac{1}{2\pi }\int_{\mathcal A} \left(\frac{\|\bm x'\|}{\|\bm x\|}\langle \bm x, \bm w^*\rangle^2 + \frac{\|\bm x\|}{\|\bm x'\|}\langle \bm x', \bm w^*\rangle^2\right)f(\bm x')\mathrm d\rho_{\mathcal X}(\bm x') + E_1^{(2)}(\bm x) + E_2^{(2)}(\bm x).
    \]
    
    Following the same argument as before we can write the integral over $\mathcal A$ as the integral over the whole space plus a term $E_3^{(2)}(\bm x)$ such that 
    \[
        |E_3^{(2)}(\bm x)| \leq \mathcal O(de^{-d \epsilon^2})\left(\|\bm x\| + \frac{\langle \bm x, \bm w^*\rangle ^2}{\|\bm x\|}\right)
    \]
    almost everywhere, and thus $\|E^{(3)}\|_{L^2(\rho_{\mathcal X})} =  \mathcal O(d^{3/2}e^{-d \epsilon^2})$.
    
    
    Therefore, we can write
    \begin{align*}
        T_{(2*)}f(\bm x) &= \frac{1}{2\pi}\int \left(\frac{\|\bm x'\|}{\|\bm x\|}\langle \bm x, \bm w^*\rangle^2 + \frac{\|\bm x\|}{\|\bm x'\|}\langle \bm x', \bm w^*\rangle^2\right)f(\bm x')\mathrm d\rho_{\mathcal X}(\bm x') + E_1^{(2)}(\bm x) + E_2^{(2)}(\bm x) + E_3^{(2)}(\bm x)\\
        & = \frac{1}{2\pi}\frac{\langle \bm x, \bm w^*\rangle^2}{\|\bm x\|} \langle \|\cdot\|, f\rangle + \frac{1}{2\pi}\|\bm x \| \left\langle \frac{\langle \cdot, \bm w^*\rangle^2}{\|\cdot\|}, f \right\rangle + E_1^{(2)}(\bm x) + E_2^{(2)}(\bm x) + E_3^{(2)}(\bm x)\\
        &:= T_{S}^{(2*)}(\bm x) + E_1^{(2)}(\bm x) + E_2^{(2)}(\bm x) + E_3^{(2)}(\bm x).
    \end{align*}

    Finally, since we chose $\epsilon = d^{-1/3}$, we have
    \begin{align*}
        \|T_{(2*)}f - T_{S}^{(2*)}f\|_{L^2(\rho_{\mathcal X})} &\leq \left(\|E_1^{(2)}\|_{L^2(\rho_{\mathcal X})} + \|E_2^{(2)}\|_{L^2(\rho_{\mathcal X})} + \|E_3^{(2)}\|_{L^2(\rho_{\mathcal X})}\right) \\
        &\leq \left(\mathcal O(d^{3/2}e^{-d^{1/3}})+ \mathcal O(d^{-1/6}) + \mathcal O(d^{3/2}e^{-d^{1/3}})\right)\\
        &= o_d(1).
    \end{align*}
\end{proof}

\subsection{Proof of Theorem \ref{thm:lin-eigenfn}}\label{app:thm-lin-eigenfn}

\begin{proof}


    We show the linear eigenfunction $\psi = \langle \cdot ,\bm v\rangle$, $\bm v \in \mathbb R^d$, continues to be an eigenfunction of $T_1$ if $\bm v = \bm w^*$ or $\bm v \perp \bm w^*$. 
    
    We use Fubini's Theorem to write
    \[
        T_1 \psi(\bm x) = \int k_1(\bm x, \bm x')\psi(\bm x')\mathrm d \rho_{\mathcal X}(\bm x') = \mathbb E_{\bm w \sim \mathcal N(0, \frac{1}{d} \bm \Gamma)}[\sigma(\langle \bm w, \bm x\rangle) \mathbb E_{\bm x' \sim \mathcal N(0, \bm I_d)}[\sigma(\langle \bm w, \bm x'\rangle)\psi(\bm x')]]
    \]
    and analytically solve these integrals. Looking at the inner expectation, and applying Stein's Lemma, we have
    \[
        \mathbb E_{\bm x' \sim \mathcal N(0, \bm I_d)}[\sigma(\langle \bm w, \bm x'\rangle)\psi(\bm x')] = \mathbb E_{\bm x' \sim \mathcal N(0, \bm I_d)}[\sigma(\langle \bm w, \bm x'\rangle)\langle \bm x', \bm v\rangle] = \mathbb E_{\bm x' \sim \mathcal N(0, \bm I_d)}[\langle D_{\bm x'}\left\{ \sigma(\langle \bm w, \bm x'\rangle) \right\}, \bm v\rangle)].
    \]

    Now, since $\sigma'(t) = \bm 1_{\{t \geq 0\}}$ almost everywhere, we have
    \[
        \mathbb E_{\bm x' \sim \mathcal N(0, \bm I_d)}[D_{\bm x'} \left\{ \sigma(\langle \bm w, \bm x'\rangle) \right\}] = \mathbb E_{\bm x' \sim \mathcal N(0, \bm I_d)}[\langle \bm w, \bm v\rangle \bm 1_{\{\langle \bm w, \bm x'\rangle \geq 0\}}] = \frac{1}{2}\langle \bm w, \bm v\rangle.
    \]

    Going back to the outer integral
    \[
        T_1\psi(\bm x) = \frac{1}{2}\mathbb E_{\bm w \sim \mathcal N(0, \frac{1}{d}\bm \Gamma)}[\sigma(\langle \bm w, \bm x\rangle) \langle \bm w, \bm v\rangle],
    \]
    and using Stein's Lemma and the same argument again
    \[
        \mathbb E_{\bm w \sim \mathcal N(0, \frac{1}{d}\bm \Gamma)}[\sigma(\langle \bm w, \bm x\rangle) \langle \bm w, \bm v\rangle] =  \frac{1}{d}\mathbb E_{\bm w \sim \mathcal N(0, \frac{1}{d}\bm \Gamma)}[\langle\bm \Gamma \nabla_{\bm w}\sigma(\langle \bm w, \bm x\rangle), \bm v\rangle] = \frac{1}{2{d}}\langle \bm \Gamma \bm x, \bm v\rangle, 
    \]
    therefore
    \[
        T_1\psi(\bm x) = \frac{1}{4d}\langle \bm \Gamma \bm x, \bm v\rangle.
    \]

    Note that the same calculation shows that
    \[
        T_0 \psi(\bm x) = \frac{1}{4d}\langle \bm x, \bm v\rangle.
    \]
    
    As a consequence, if we let $\bm v = \bm w^*$, we have
    \[
        \frac{1}{4d}\langle \bm \Gamma \bm x, \bm v\rangle = \frac{1}{4d}\langle \bm \Gamma \bm x, \bm w^*\rangle = \left(\frac{A}{4d} + \frac{B}{4d}\right)\langle\bm x, \bm w^*\rangle,
    \]
        
    And, on the other hand, if $\bm v \perp \bm w^*$
    \[
        \frac{1}{4d}\langle \bm \Gamma \bm x, \bm v\rangle = \frac{A}{4d}\langle\bm x, \bm v\rangle.        
    \]

\end{proof}

\subsection{Proof of Theorem \ref{thm:approx-top-eigenfn}}\label{app:thm-approx-top-eigenfn}
\begin{proof}

    We have from Lemma \ref{lemma:approx-first-order-op} that, for a given $f$ that can be represented by the eigenfunctions of the kernel, the action of operator $S$ can be approximated by
    \[
        T_Sf = T_S^{{(1*)}}f + T_S^{{(2*)}}f + E.
    \]
    If we recall the closed form of the first term
    \[
        T_S^{{(1*)}}f(\bm x) = \frac{\langle \bm x, \bm w^*\rangle}{2}\langle\langle \cdot, \bm w^*\rangle,f\rangle, 
    \]
    we note the inner product term implies that any function that is orthogonal to the linear function must not interact with this term. 
    
    Also note that Theorem \ref{thm:lin-eigenfn} shows the linear function is a true eigenfunction of $T_1$, thus if a different function happens to be the top eigenfunction of the operator, then they must be orthogonal to each other. 
    
    Therefore, to understand the action
    \[
        T_1 f(\bm x) = AT_0f(\bm x) + \frac{B}{2d}T_S f(\bm x) + T_Rf(\bm x)
    \]
    on any function other than the linear function it suffices to solve the eigenvalue problem for $AT_0 + \frac{B}{2d}T_{S}^{(2*)}$.

    {\bf Solving the eigenvalue problem for $AT_0 + \frac{B}{2d}T_{S}^{(2*)}$:}
    Recall that
    \[
        T_{S}^{(2*)}f (\bm x) = \frac{1}{2\pi}\left( \|\bm x \| \left\langle \frac{\langle \cdot, \bm w^*\rangle^2}{\|\cdot\|}, f \right\rangle + \frac{\langle \bm x, \bm w^*\rangle^2}{\|\bm x\|} \langle \|\cdot\|, f\rangle \right)
    \]
    and if we define $f_1(\bm x) := \|\bm x\|$ and $f_2(\bm x) := \frac{\langle \bm x, \bm w^*\rangle^2}{\|\bm x\|}$ we have
    \[
        T_{S}^{(2*)}f = \frac{1}{2\pi}\Big( f_1\langle f_2, f\rangle + f_2 \langle f_1, f\rangle\Big)\, ,
    \]
    therefore the eigenfunctions must lie in $\text{span}\{f_1, f_2\}$. To bridge the gap between both operators, we recall the eigenbasis of the isotropic operator given by Lemma \ref{lemma:iso-eigenfn}, and we translate our functions $f_1$ and $f_2$ to that basis.
        
    First, we let $r = \|\bm x\|$ and $\bm \omega = \frac{\bm x}{\|\bm x\|}$ and consider the following elements of the spherical harmonics basis
    \[
        Y_0(\bm \omega) = 1, \quad \hat{Y}_2(\bm \omega, \bm w^*) = \langle \bm \omega, \bm w^*\rangle^2 - \frac{1}{d}.
    \]
    Note the notation $\hat{Y}_2$ to indicate this function is not normalized in $L^2(\rho_{\mathcal X})$.
    With these we can write
    \[
        f_1(\bm x) = \|\bm x\| = r Y_0(\bm \omega), \quad f_2(\bm x) = \frac{\langle \bm x, \bm w^*\rangle^2}{\|\bm x\|} = r\langle \bm \omega, \bm w^*\rangle^2,
    \]
    and writing the inner product as a combination of both functions yields
    \[
        \langle \bm \omega, \bm w^*\rangle^2 = \hat{Y}_2(\bm \omega, \bm w^*) + \frac{1}{d}Y_0(\bm \omega),
    \]
    \[
        f_1(\bm x) = r Y_0(\bm \omega), \quad f_2(\bm x) = r \hat{Y}_2(\bm \omega, \bm w^*) + \frac{r}{d} Y_0(\bm \omega).
    \]
    
    Next, we solve the eigenvalue problem of the combined operator $AT_0 +\frac{B}{2d} T_{S}^{(2*)}$. Let $\varphi$ be a non trivial eigenfunction of this operator. Since $T_0$ is fully diagonalized by $r Y_0$ and $r \hat{Y}_2$ and since the action of $T_{S}^{(2*)}$ is restricted to $\text{span}\{r Y_0, r \hat{Y}_2\}$, we know $\varphi$ must take the form
    \[
        \varphi = c_0 [r Y_0(\bm \omega)] + c_2[r \hat{Y}_2(\bm \omega, \bm w^*)]
    \]
    for some coefficients $c_0$ and $c_2$.
    
    If we apply
    \[
        T_{S}^{(2*)}\varphi = \frac{\langle f_2, \varphi \rangle}{2 \pi}f_1 + \frac{ \langle f_1, \varphi \rangle}{2 \pi}f_2,
    \]
    using the identities $f_1 = r Y_0(\omega)$ and $f_2 = r \hat{Y}_2(\omega, \bm w^*) + \frac{r}{d}Y_0(\omega)$, we obtain
    \begin{align*}
        T_{S}^{(2*)}\varphi &= \left\langle r \hat{Y}_2 + \frac{r}{d} Y_0, c_0[r Y_0] + c_2[r \hat{Y}_2]\right\rangle \frac{r Y_0}{2 \pi}  + \left\langle r Y_0,  c_0[r Y_0] + c_2[r \hat{Y}_2]\right\rangle\frac{1}{2 \pi}\left( r \hat{Y}_2 + \frac{r}{d} Y_0 \right)\\
        &= \frac{r Y_0}{2 \pi} \left( \frac{c_0}{d} \langle r Y_0, r Y_0\rangle + c_2 \langle r \hat{Y}_2, r \hat{Y}_2\rangle\right) + \frac{1}{ 2 \pi}\left( r \hat{Y}_2 + \frac{r}{d} Y_0 \right)c_0\langle r Y_0, r Y_0\rangle\\
        &= \frac{r Y_0}{2 \pi}\left( \frac{2c_0}{d} \langle r Y_0, r Y_0\rangle + c_2 \langle r \hat{Y}_2, r \hat{Y}_2\rangle\right) + \frac{r \hat{Y}_2}{2 \pi} \Big(c_0\langle r Y_0, r Y_0\rangle \Big).
    \end{align*}

    This translates to the matrix operator on the vector of coefficients
    \[
        T_{S}^{(2*)}\varphi = \frac{1}{2\pi}\begin{bmatrix}
            \frac{2}{d}\langle r Y_0, r Y_0\rangle & \langle r \hat{Y}_2, r \hat{Y}_2\rangle\\
            \langle r Y_0, r Y_0\rangle & 0
        \end{bmatrix} \begin{pmatrix}
            c_0\\
            c_2
        \end{pmatrix} \cdot \begin{pmatrix}
            r Y_0\\
            r \hat{Y}_2
        \end{pmatrix}
    \]
    and since $\langle r Y_0, r Y_0\rangle = d$ and $\langle r \hat{Y}_2,  r \hat{Y}_2\rangle = \frac{2(d-1)}{d(d+2)}$ we have
    \[
        T_{S}^{(2*)}\varphi = \frac{1}{2\pi}\begin{bmatrix}
            2 & \frac{2(d-1)}{d(d+2)}\\
            d & 0
        \end{bmatrix} \begin{pmatrix}
            c_0\\
            c_2
        \end{pmatrix} \cdot \begin{pmatrix}
            r Y_0\\
            r \hat{Y}_2
        \end{pmatrix} .
    \]

    Furthermore, because these are non normalized orthogonal eigenfunctions of $T_0$, they perfectly diagonalize the operator. Thus, if $\lambda_{\max}(T_0)$ and $\lambda_2(T_0)$ are the eigenvalues for $r Y_0$ and $rY_2$ respectively, we have the following
    \[
        AT_0 \varphi = \begin{bmatrix}
            A\lambda_{\max}(T_0) & 0\\
            0 & A\lambda_2(T_0)
        \end{bmatrix} \begin{pmatrix}
            c_0\\
            c_2
        \end{pmatrix} \cdot \begin{pmatrix}
            r Y_0\\
            r \hat{Y}_2
        \end{pmatrix},
    \]
    and the combined action is simply the combination of these matrices
    \[
        \left( AT_0 + \frac{B}{2d}T_{S}^{(2*)} \right) \varphi = \begin{bmatrix}
            A\lambda_{\max}(T_0) + \frac{B}{2\pi d} & \frac{B(d-1)}{2 \pi d^2(d+2)}\\
            \frac{B}{4 \pi } & A\lambda_2(T_0)
        \end{bmatrix} \begin{pmatrix}
            c_0\\
            c_2
        \end{pmatrix} \cdot \begin{pmatrix}
            r Y_0\\
            r \hat{Y}_2
        \end{pmatrix} .
    \]
    Note that $\lambda_2$ represents the eigenvalue corresponding to a single spherical harmonic of degree 2. We know the dimension of the degree-2 subspace on $\mathbb {S}^{d-1}$ is given by the degeneracy formula $N(d, 2) = \frac{(d-1)(d+2)}{2}$, thus since the total macroscopic energy is distributed uniformly across the orthogonal basis in the subspace, the individual eigenvalue scales as $\lambda_2 = \mathcal{O}(d^{-2})$. 

    We draw attention to the $\mathcal O(B)$ coefficient on the lower left entry of the matrix which implies a strong coupling between these functions under the combined operator.
    
    Again, solving for the roots of the characteristic polynomial, we find the following two eigenvalues
    \[
        \tilde{\lambda}_{\pm} = \frac{1}{2} \left[ \left( A \lambda_{\max}(T_0) + \frac{B}{2 \pi d} + A \lambda_2(T_0) \right) \pm \sqrt{ \left( A \lambda_{\max}(T_0) + \frac{B}{2 \pi d} - A \lambda_2(T_0) \right)^2 + \frac{B^2(d-1)}{2 \pi^2 d^2(d+2)} } \right],
    \]
    or, more intuitively, using the Taylor expansion of $\sqrt{x + \epsilon}$ we can write
    \[
        \tilde{\lambda}_+ = A\lambda_{\max}(T_0) + \frac{B}{2 \pi d} + o_d\left(\frac{|B|}{d}\right), \quad \tilde{\lambda}_- = \mathcal O\left( A\lambda_2(T_0) + \frac{B^2}{d^3}\right).
    \]
    
    Furthermore, we know that any eigenfunction with this eigenvalue must satisfy
    \[
         \begin{bmatrix}
            A\lambda_{\max}(T_0) + \frac{B}{2\pi d} & \frac{B(d-1)}{2 \pi d^2(d+2)}\\
            \frac{B}{4 \pi } & A\lambda_2(T_0)
        \end{bmatrix}\begin{pmatrix}
            c_0\\
            c_2
        \end{pmatrix} = \lambda_{\pm}\begin{pmatrix}
            c_0\\
            c_2
        \end{pmatrix}
    \]
    which implies the relationship
    \[
        \frac{B}{4\pi}c_0 = (\lambda_{\pm} - A \lambda_2(T_0))c_2.
    \]
    
    Since the eigenfunctions can be arbitrarily scaled, we choose $c_0 = \frac{1}{\sqrt d}$ to reconstruct the isotropic eigenfunction $c_0 \|\bm x\| Y_0 = \frac{\|\bm x\|}{\sqrt d} Y_0$, and solving for $c_2$ gives
    \[
        c_2 = \frac{B}{4 \pi \sqrt d(\tilde{\lambda}_{\pm}-A\lambda_2(T_0))}
    \]
    which lets us define the new eigenfunctions as
    \[
        \hat{\Psi}_{\pm}(\bm x) = \frac{\|\bm x\|}{\sqrt d} Y_0(\bm \omega)+ \left[ \frac{B}{4 \pi( \tilde{\lambda}_{\pm} - A\lambda_2(T_0) )} \right] \frac{\|\bm x\|}{\sqrt d}\hat{Y}_2(\bm \omega, \bm w^*) \, .
    \]

    To ensure this function is normalized, we calculate
    \[
        \|\hat{\Psi}_{\pm}\|^2_{L^2(\rho_{\mathcal X})} = 1 + \frac{B^2}{16 \pi ^2( \tilde{\lambda}_{\pm}-A\lambda_2(T_0))^2} \frac{2d -2}{d^2(d+2)}:= N_{\pm}
    \]
    and therefore the normalized approximate eigenfunctions are given by
    \[
        \tilde{\Psi}_{\pm}(\bm x) = \frac{1}{\sqrt{N_{\pm}}}\left( \frac{\|\bm x\|}{\sqrt d} Y_0(\bm \omega)+ \left[ \frac{B}{4 \pi ( \tilde{\lambda}_{\pm}-A\lambda_2(T_0) )} \right] \frac{\|\bm x\|}{\sqrt d} \hat{Y}_2(\bm \omega, \bm w^*)\right)\, .
    \]

    Lastly, to explicitly show the magnitude of the quadratic feature, if $Y_2$ is the normalized zonal harmonics, we write $\hat{Y}_2 = \|\hat{Y}_2\|_{L^2(\rho_{\mathcal X})} Y_2$ with $\|\hat{Y}_2\|_{L^2(\rho_{\mathcal X})} = \frac{1}{d}\sqrt{\frac{2d - 2}{d+2}}$. Therefore, we have the final form for the eigenfunctions:
    \[
        \tilde{\Psi}_{\pm}(\bm x) = \frac{1}{\sqrt{N_{\pm}}}\frac{\|\bm x\|}{\sqrt d} \Big[ Y_0(\bm \omega)+  \tau_{\pm} Y_2(\bm \omega, \bm w^*) \Big]\,,
    \]
    where we define the alignment magnitudes by the quantity
    \[
        \tau_{\pm} := \frac{1}{d}\sqrt{\frac{2d - 2}{d+2}}\left[\frac{B}{4 \pi (\tilde{\lambda}_{\pm}-A\lambda_2(T_0))}\right].
    \]

    From now on we study the approximate eigenfunction $\tilde{\Psi}_{+}$ with associated approximate eigenvalue $\lambda_{+}$, since we clearly have $\tilde{\lambda}_{+} > \tilde{\lambda}_{-}$.
    
    {\bf The approximate eigenvalue $\tilde{\lambda}_{+}$ dominates every other eigenvalue of $T_1$:}
    From the previous derivations, if we look at the approximate eigenvalue as functions of the dimension $\lambda_{+}(d)$, we established that 
    \[
        \lim_{d \to \infty} \tilde{\lambda}_{+}(d) = A\lambda_{\max}(T_0) > 0
    \]
    Thus, for any chosen $\epsilon > 0$, there exists an integer $d_0$ such that for all $d > d_0$
    \[
        \tilde{\lambda}_{+}(d) > A\lambda_{\max}(T_0) - \epsilon
    \]
    Choose $\epsilon = \frac{A\lambda_{\max}(T_0)}{2}$ such that there exists a $d_0$ such that for all $d > d_0$
    \[
        \tilde{\lambda}_{+}(d) > \frac{A\lambda_{\max}(T_0)}{2}.
    \]
    
    By Theorem \ref{thm:lin-eigenfn}, we know $\lambda_1(T_1) = A\lambda_1(T_0) + \frac{B}{4d} = \mathcal{O}(d^{-1})$. Thus, there exists a real constant $M > 0$ and an integer $d_1$ such that for all $d > d_1$
    \[
        \lambda_1(T_1) \le \frac{M}{d}.
    \]
    Because the ReLU kernel has monotonically decreasing eigenvalues, the linear eigenvalue serves as a strict upper bound for all higher-degree subspaces
    \[
        \lambda_1(T_0) > \lambda_k(T_0) \quad \text{for all } k \ge 2.
    \]
    Furthermore, by Theorem \ref{thm:eigenv-decay} there exists a radius $R > 0$ and $\varepsilon > 0$ such that for constant $C_R > 0$ independent of the dimension we have
    \[
        \lambda_k(T_1) \leq C_R\lambda_k(T_0) + \varepsilon < 2C_R\lambda_1(T_0)
    \]
    for all $k \geq 1$ which implies
    \[
        \lambda_k (T_1) \leq \frac{2C_RM}{d}\, .
    \]
    
    Thus, if we can prove $\tilde{\lambda}_{+}(d) > \lambda_1(T_1)$, we automatically prove it dominates all $\lambda_k(T_1)$ for $k \ge 1$.
    
    We want to guarantee that $\tilde{\lambda}_{+}(d) > \lambda_1(T_1)$ which is true as long as
    \[
        \frac{A\lambda_{\max}(T_0)}{2} > \frac{2C_RM}{d} \implies d > \frac{4C_RM}{A\lambda_{\max}(T_0)}.
    \]
    
    Define 
    \[
        d^* = \max\left(d_0, d_1, \left\lfloor \frac{4C_RM}{A\lambda_{\max}(T_0)} \right\rfloor + 1 \right)
    \]
    then for any dimension $d > d^*$, the following holds
    \[
        \tilde{\lambda}_{+}(d) > \frac{A\lambda_{\max}(T_0)}{2} > \frac{2C_RM}{d} \ge \lambda_1(T_1) \geq  \lambda_{k \ge 2}(T_1).
    \]

    {\bf The approximate eigenfunction and the true eigenfunction are close in norm:}
    Finally, to show $\tilde{\Psi}_{+}$ is close to the original top eigenfunction, we denote by $\Psi$ the top eigenfunction and use the expansion to write
    \[
        T_1 \tilde{\Psi}_{+} = \left(AT_0 + \frac{B}{2d}T_{S_{(1*)}} + \frac{B}{2d}T_{S}^{(2*)}\right)\tilde{\Psi} + e = \tilde{\lambda}_{+}\tilde\Psi_{+} + e
    \]
    with $\| e \| = o_d\left(\frac{|B|}{d}\right)$ and we write this as
    \[
        (T_1 - \tilde{\lambda}_{+} I)\tilde{\Psi}_{+} = e\, .
    \]
    Now, consider the projection onto the top eigenspace of $T_1$ defined by
    \[
        P_{\text{top}}f =  \langle \Psi, f\rangle \Psi,
    \]
    and its orthogonal complement $P_{\perp} = I - P_{\text{top}}$. Expanding the action of the combined operator using these projections we get
    \[
        (T_1 - \tilde{\lambda}_{+} I)\tilde{\Psi}_{+} = (T_1 - \tilde{\lambda}_{+} I)(P_{\text{top}}\tilde{\Psi}_{+} + P_{\perp}\tilde{\Psi}_{+}) = e
    \]
    and since $T_1(P_{\text{top}} f) = \lambda_{\max}(T_1) P_{\text{top}} f$ this simplifies to
    \[
        [\lambda_{\max}(T_1) - \tilde{\lambda}_{+}]P_{\text{top}} \tilde{\Psi}_{+} + (T_1 - \tilde{\lambda}_{+} I)P_{\perp}\tilde{\Psi}_{+} = e.
    \]
    Since we have only orthogonal objects, the norm of this expression is equal to
    \[
        |\lambda_{\max}(T_1)- \lambda_{+}|^2\| P_{\text{top}} \tilde{\Psi}_{+} \|_{L^2(\rho_{\mathcal X})} ^2 + \|(T_1 - \tilde{\lambda}_{+} I)P_{\perp}\tilde{\Psi}_{+}\|_{L^2(\rho_{\mathcal X})} ^2 = o_d\left(\frac{B^2}{d^2}\right)
    \]
    which immediately implies
    \[
        \|(T_1 - \tilde{\lambda}_{+} I)P_{\perp}\tilde{\Psi}_{+}\|_{L^2(\rho_{\mathcal X})} ^2 = o_d\left(\frac{B^2}{d^2}\right)\, ,
    \]
    and we will use this to bound $\| P_{\perp} \tilde{\Psi}\|$. If we define the approximate spectral gap by $\tilde{\delta} = \inf_{\mu \in \sigma(T_1)\setminus\{\lambda_{\max}(T_1)\}}|\mu - \tilde{\lambda}_{+}|$, we have
    \[
        \| (T_1 - \tilde{\lambda}_{+} I)P_{\perp}\tilde{\Psi}_{+} \|_{L^2(\rho_{\mathcal X})}  \geq  \tilde{\delta}\| P_{\perp}\tilde{\Psi}_{+} \|_{L^2(\rho_{\mathcal X})} 
    \]
    and
    \[
        \tilde{\delta} \| P_{\perp}\tilde{\Psi}_{+} \|_{L^2(\rho_{\mathcal X})}  = o_d\left(\frac{|B|}{d}\right).
    \]

    From Theorem \ref{thm:eigenv-decay}, we have the lower bound
    \[
        c\lambda_{\max}(T_0) \leq \lambda_{\max}(T_1)
    \]
    where $c > 0$ is an absolute constant. We know $\lambda_{\max}(T_0) = \Theta(1)$ thus $c\lambda_{\max}(T_0)$ is bounded away from zero. which gives the lower bound
    \[
        \lambda_{\max}(T_1) \geq c\lambda_{\max}(T_0) = \Omega(1)\, .
    \]

    Now if $\delta := \lambda_{\max}(T_1) - \lambda_1(T_1)$ is the true spectral gap, we know from Theorem \ref{thm:lin-eigenfn} that $\lambda_1(T_1) = \Theta(1/d)$. Substituting our lower bound for $\lambda_{\max}(T_1)$ into the gap we have
    \[
        \delta \geq c\lambda_{\max}(T_0) - \lambda_1(T_1) = \Omega(1)\, ,
    \]
    for a high enough dimension $d$. The difference between the gaps is given by
    \[
        |\lambda_{\max}(T_1) - \tilde{\lambda}_{+}| = |\delta - \tilde{\delta}|
    \]
    and since $|\lambda_{\max}(T_1) - \tilde{\lambda}_{+}| = o_d(|B|/d)$ the triangle inequality implies
    \[
        \tilde{\delta} \geq \delta - o_d\left(\frac{|B|}{d}\right) \, ,
    \]
    and we know $\tilde{\delta}$ must bounded away from zero independently of $d$. Thus, going back to bounding the norm of $P_\perp \tilde{\Psi}_{+}$, we obtain
    \[
         \|P_{\perp}\tilde{\Psi}_{+} \|_{L^2(\rho_{\mathcal X})}   = \frac{1}{\tilde{\delta}} o_d\left(\frac{|B|}{d}\right) = o_d\left(\frac{|B|}{d}\right)\, .
    \]
    
    To bound the difference under the norm, we note that $\|\tilde{\Psi}\| = 1$ and write
    \[
        \| \tilde{\Psi}_{+} \|_{L^2(\rho_{\mathcal X})} ^2 = |\langle \tilde{\Psi}_{+}, \Psi\rangle|^2 + \|P_{\perp} \tilde{\Psi}_{+}\|_{L^2(\rho_{\mathcal X})} ^2
    \]
    \[
        |\langle \tilde{\Psi}_{+}, \Psi\rangle|^2 = 1 - o_d\left(\frac{B^2}{d^2}\right)
    \]
    Finally, combining all together we have
    \begin{align*}
        \| \tilde{\Psi}_{+} - \Psi\|_{L^2(\rho_{\mathcal X})} ^2 &= \| \tilde{\Psi}_{+} \|_{L^2(\rho_{\mathcal X})} ^2 -2\langle \tilde{\Psi}_{+}, \Psi\rangle + \|\Psi\|_{L^2(\rho_{\mathcal X})} ^2\\
        &= 2 - 2\left[1- o_d\left(\frac{B^2}{d^2}\right)\right]\\
        &=  o_d\left(\frac{B^2}{d^2}\right)
    \end{align*}
    and the proof is completed.
\end{proof}

